% Ngân hàng đề Toán 12 (KNTT) -- Chương II, Bài 8, Dạng 10
% Dạng: Tọa độ hóa hình học: tam giác, hình chóp, điểm cách đều
%% Phương án ĐẦU TIÊN trong \cauTN luôn là đáp án đúng (gói trondethi tự đảo).

\cauTN[id=C2-B8-D10-NB-001,muc=NB,loigiai={Ta có $\overrightarrow{AB}=(1;-3;-2)$, $\overrightarrow{AC}=(-5;-1;-3)$. Diện tích $S=\dfrac{1}{2}\sqrt{|\overrightarrow{AB}|^2|\overrightarrow{AC}|^2-(\overrightarrow{AB}\cdot\overrightarrow{AC})^2}=\dfrac{1}{2}\sqrt{14\cdot35-4^2}=\dfrac{\sqrt{474}}{2}$.}]{Trong không gian $Oxyz$, cho $A(2;1;3)$, $B(3;-2;1)$, $C(-3;0;0)$. Diện tích tam giác $ABC$ bằng}
  {$\dfrac{\sqrt{474}}{2}$.}
  {$\sqrt{474}$.}{$\sqrt{237}$.}{$\dfrac{\sqrt{478}}{2}$.}

\cauTN[id=C2-B8-D10-NB-002,muc=NB,loigiai={Ta có $\overrightarrow{AB}=(-2;-5;-2)$, $\overrightarrow{AC}=(-2;-1;2)$. Diện tích $S=\dfrac{1}{2}\sqrt{|\overrightarrow{AB}|^2|\overrightarrow{AC}|^2-(\overrightarrow{AB}\cdot\overrightarrow{AC})^2}=\dfrac{1}{2}\sqrt{33\cdot9-5^2}=2\sqrt{17}$.}]{Trong không gian $Oxyz$, cho $A(3;2;-1)$, $B(1;-3;-3)$, $C(1;1;1)$. Diện tích tam giác $ABC$ bằng}
  {$2\sqrt{17}$.}
  {$4\sqrt{17}$.}{$2\sqrt{34}$.}{$\sqrt{69}$.}

\cauTN[id=C2-B8-D10-NB-003,muc=NB,loigiai={Ta có $\overrightarrow{AB}=(4;0;-6)$, $\overrightarrow{AC}=(2;2;-5)$. Diện tích $S=\dfrac{1}{2}\sqrt{|\overrightarrow{AB}|^2|\overrightarrow{AC}|^2-(\overrightarrow{AB}\cdot\overrightarrow{AC})^2}=\dfrac{1}{2}\sqrt{52\cdot33-38^2}=2\sqrt{17}$.}]{Trong không gian $Oxyz$, cho $A(-3;1;3)$, $B(1;1;-3)$, $C(-1;3;-2)$. Diện tích tam giác $ABC$ bằng}
  {$2\sqrt{17}$.}
  {$4\sqrt{17}$.}{$2\sqrt{34}$.}{$\sqrt{69}$.}

\cauTN[id=C2-B8-D10-NB-004,muc=NB,loigiai={Ta có $\overrightarrow{AB}=(0;2;0)$, $\overrightarrow{AC}=(5;3;5)$. Diện tích $S=\dfrac{1}{2}\sqrt{|\overrightarrow{AB}|^2|\overrightarrow{AC}|^2-(\overrightarrow{AB}\cdot\overrightarrow{AC})^2}=\dfrac{1}{2}\sqrt{4\cdot59-6^2}=5\sqrt{2}$.}]{Trong không gian $Oxyz$, cho $A(-3;-1;-2)$, $B(-3;1;-2)$, $C(2;2;3)$. Diện tích tam giác $ABC$ bằng}
  {$5\sqrt{2}$.}
  {$10\sqrt{2}$.}{$10$.}{$\sqrt{51}$.}

\cauTN[id=C2-B8-D10-NB-005,muc=NB,loigiai={Ta có $\overrightarrow{AB}=(0;2;-5)$, $\overrightarrow{AC}=(1;3;1)$. Diện tích $S=\dfrac{1}{2}\sqrt{|\overrightarrow{AB}|^2|\overrightarrow{AC}|^2-(\overrightarrow{AB}\cdot\overrightarrow{AC})^2}=\dfrac{1}{2}\sqrt{29\cdot11-1^2}=\dfrac{\sqrt{318}}{2}$.}]{Trong không gian $Oxyz$, cho $A(-3;-1;2)$, $B(-3;1;-3)$, $C(-2;2;3)$. Diện tích tam giác $ABC$ bằng}
  {$\dfrac{\sqrt{318}}{2}$.}
  {$\sqrt{318}$.}{$\sqrt{159}$.}{$\dfrac{\sqrt{322}}{2}$.}

\cauTN[id=C2-B8-D10-NB-006,muc=NB,loigiai={Ta có $\overrightarrow{AB}=(-2;2;-1)$, $\overrightarrow{AC}=(-4;0;1)$. Diện tích $S=\dfrac{1}{2}\sqrt{|\overrightarrow{AB}|^2|\overrightarrow{AC}|^2-(\overrightarrow{AB}\cdot\overrightarrow{AC})^2}=\dfrac{1}{2}\sqrt{9\cdot17-7^2}=\sqrt{26}$.}]{Trong không gian $Oxyz$, cho $A(1;-2;-1)$, $B(-1;0;-2)$, $C(-3;-2;0)$. Diện tích tam giác $ABC$ bằng}
  {$\sqrt{26}$.}
  {$2\sqrt{26}$.}{$2\sqrt{13}$.}{$3\sqrt{3}$.}

\cauTN[id=C2-B8-D10-NB-007,muc=NB,loigiai={Ta có $\overrightarrow{AB}=(4;5;3)$, $\overrightarrow{AC}=(4;3;1)$. Diện tích $S=\dfrac{1}{2}\sqrt{|\overrightarrow{AB}|^2|\overrightarrow{AC}|^2-(\overrightarrow{AB}\cdot\overrightarrow{AC})^2}=\dfrac{1}{2}\sqrt{50\cdot26-34^2}=6$.}]{Trong không gian $Oxyz$, cho $A(-2;-3;-2)$, $B(2;2;1)$, $C(2;0;-1)$. Diện tích tam giác $ABC$ bằng}
  {$6$.}
  {$12$.}{$6\sqrt{2}$.}{$\sqrt{37}$.}

\cauTN[id=C2-B8-D10-NB-008,muc=NB,loigiai={Ta có $\overrightarrow{AB}=(6;1;-2)$, $\overrightarrow{AC}=(0;-3;-4)$. Diện tích $S=\dfrac{1}{2}\sqrt{|\overrightarrow{AB}|^2|\overrightarrow{AC}|^2-(\overrightarrow{AB}\cdot\overrightarrow{AC})^2}=\dfrac{1}{2}\sqrt{41\cdot25-5^2}=5\sqrt{10}$.}]{Trong không gian $Oxyz$, cho $A(-3;2;1)$, $B(3;3;-1)$, $C(-3;-1;-3)$. Diện tích tam giác $ABC$ bằng}
  {$5\sqrt{10}$.}
  {$10\sqrt{10}$.}{$10\sqrt{5}$.}{$\sqrt{251}$.}

\cauTN[id=C2-B8-D10-NB-009,muc=NB,loigiai={Ta có $\overrightarrow{AB}=(-1;-1;5)$, $\overrightarrow{AC}=(-3;1;6)$. Diện tích $S=\dfrac{1}{2}\sqrt{|\overrightarrow{AB}|^2|\overrightarrow{AC}|^2-(\overrightarrow{AB}\cdot\overrightarrow{AC})^2}=\dfrac{1}{2}\sqrt{27\cdot46-32^2}=\dfrac{\sqrt{218}}{2}$.}]{Trong không gian $Oxyz$, cho $A(1;-1;-3)$, $B(0;-2;2)$, $C(-2;0;3)$. Diện tích tam giác $ABC$ bằng}
  {$\dfrac{\sqrt{218}}{2}$.}
  {$\sqrt{218}$.}{$\sqrt{109}$.}{$\dfrac{\sqrt{222}}{2}$.}

\cauTN[id=C2-B8-D10-NB-010,muc=NB,loigiai={Ta có $\overrightarrow{AB}=(-2;4;3)$, $\overrightarrow{AC}=(3;2;3)$. Diện tích $S=\dfrac{1}{2}\sqrt{|\overrightarrow{AB}|^2|\overrightarrow{AC}|^2-(\overrightarrow{AB}\cdot\overrightarrow{AC})^2}=\dfrac{1}{2}\sqrt{29\cdot22-11^2}=\dfrac{\sqrt{517}}{2}$.}]{Trong không gian $Oxyz$, cho $A(0;-1;0)$, $B(-2;3;3)$, $C(3;1;3)$. Diện tích tam giác $ABC$ bằng}
  {$\dfrac{\sqrt{517}}{2}$.}
  {$\sqrt{517}$.}{$\dfrac{\sqrt{1034}}{2}$.}{$\dfrac{\sqrt{521}}{2}$.}

\cauTN[id=C2-B8-D10-NB-011,muc=NB,loigiai={Ta có $\overrightarrow{AB}=(1;-3;1)$, $\overrightarrow{AC}=(-3;1;-4)$. Diện tích $S=\dfrac{1}{2}\sqrt{|\overrightarrow{AB}|^2|\overrightarrow{AC}|^2-(\overrightarrow{AB}\cdot\overrightarrow{AC})^2}=\dfrac{1}{2}\sqrt{11\cdot26+10^2}=\dfrac{\sqrt{186}}{2}$.}]{Trong không gian $Oxyz$, cho $A(1;0;2)$, $B(2;-3;3)$, $C(-2;1;-2)$. Diện tích tam giác $ABC$ bằng}
  {$\dfrac{\sqrt{186}}{2}$.}
  {$\sqrt{186}$.}{$\sqrt{93}$.}{$\dfrac{\sqrt{190}}{2}$.}

\cauTN[id=C2-B8-D10-NB-012,muc=NB,loigiai={Ta có $\overrightarrow{AB}=(-4;-3;-2)$, $\overrightarrow{AC}=(-4;0;-2)$. Diện tích $S=\dfrac{1}{2}\sqrt{|\overrightarrow{AB}|^2|\overrightarrow{AC}|^2-(\overrightarrow{AB}\cdot\overrightarrow{AC})^2}=\dfrac{1}{2}\sqrt{29\cdot20-20^2}=3\sqrt{5}$.}]{Trong không gian $Oxyz$, cho $A(2;1;1)$, $B(-2;-2;-1)$, $C(-2;1;-1)$. Diện tích tam giác $ABC$ bằng}
  {$3\sqrt{5}$.}
  {$6\sqrt{5}$.}{$3\sqrt{10}$.}{$\sqrt{46}$.}

\cauTN[id=C2-B8-D10-NB-013,muc=NB,loigiai={Gọi $M$ thuộc trục $Oz$, tọa độ $M$ chỉ khác $0$ ở thành phần tương ứng. Từ $MA=MB$ suy ra $MA^2=MB^2$, giải phương trình bậc nhất ta được $M=(0;0;\dfrac{7}{2})$.}]{Trong không gian $Oxyz$, cho $A(0;1;1)$, $B(-2;-1;2)$. Tìm tọa độ điểm $M$ thuộc trục $Oz$ sao cho $M$ cách đều hai điểm $A$ và $B$.}
  {$(0;0;\dfrac{7}{2})$.}
  {$(0;0;\dfrac{5}{2})$.}{$(0;1;\dfrac{7}{2})$.}{$(0;-2;\dfrac{7}{2})$.}

\cauTN[id=C2-B8-D10-NB-014,muc=NB,loigiai={Gọi $M$ thuộc trục $Oz$, tọa độ $M$ chỉ khác $0$ ở thành phần tương ứng. Từ $MA=MB$ suy ra $MA^2=MB^2$, giải phương trình bậc nhất ta được $M=(0;0;0)$.}]{Trong không gian $Oxyz$, cho $A(-3;1;-3)$, $B(3;3;1)$. Tìm tọa độ điểm $M$ thuộc trục $Oz$ sao cho $M$ cách đều hai điểm $A$ và $B$.}
  {$(0;0;0)$.}
  {$(0;1;0)$.}{$(-2;0;0)$.}{$(-1;0;0)$.}

\cauTN[id=C2-B8-D10-NB-015,muc=NB,loigiai={Gọi $M$ thuộc trục $Oy$, tọa độ $M$ chỉ khác $0$ ở thành phần tương ứng. Từ $MA=MB$ suy ra $MA^2=MB^2$, giải phương trình bậc nhất ta được $M=(0;\dfrac{3}{2};0)$.}]{Trong không gian $Oxyz$, cho $A(1;-3;1)$, $B(3;-2;1)$. Tìm tọa độ điểm $M$ thuộc trục $Oy$ sao cho $M$ cách đều hai điểm $A$ và $B$.}
  {$(0;\dfrac{3}{2};0)$.}
  {$(-2;\dfrac{3}{2};0)$.}{$(0;\dfrac{3}{2};1)$.}{$(0;\dfrac{1}{2};0)$.}

\cauTN[id=C2-B8-D10-NB-016,muc=NB,loigiai={Gọi $M$ thuộc trục $Oz$, tọa độ $M$ chỉ khác $0$ ở thành phần tương ứng. Từ $MA=MB$ suy ra $MA^2=MB^2$, giải phương trình bậc nhất ta được $M=(0;0;0)$.}]{Trong không gian $Oxyz$, cho $A(-1;3;-2)$, $B(-3;-2;-1)$. Tìm tọa độ điểm $M$ thuộc trục $Oz$ sao cho $M$ cách đều hai điểm $A$ và $B$.}
  {$(0;0;0)$.}
  {$(0;0;1)$.}{$(0;2;0)$.}{$(0;0;-2)$.}

\cauTN[id=C2-B8-D10-NB-017,muc=NB,loigiai={Gọi $M$ thuộc trục $Oz$, tọa độ $M$ chỉ khác $0$ ở thành phần tương ứng. Từ $MA=MB$ suy ra $MA^2=MB^2$, giải phương trình bậc nhất ta được $M=(0;0;-\dfrac{2}{5})$.}]{Trong không gian $Oxyz$, cho $A(0;-1;2)$, $B(0;0;-3)$. Tìm tọa độ điểm $M$ thuộc trục $Oz$ sao cho $M$ cách đều hai điểm $A$ và $B$.}
  {$(0;0;-\dfrac{2}{5})$.}
  {$(0;-\dfrac{2}{5};0)$.}{$(-\dfrac{2}{5};0;0)$.}{$(0;0;\dfrac{2}{5})$.}

\cauTN[id=C2-B8-D10-NB-018,muc=NB,loigiai={Gọi $M$ thuộc trục $Oz$, tọa độ $M$ chỉ khác $0$ ở thành phần tương ứng. Từ $MA=MB$ suy ra $MA^2=MB^2$, giải phương trình bậc nhất ta được $M=(0;0;3)$.}]{Trong không gian $Oxyz$, cho $A(-2;2;0)$, $B(3;2;1)$. Tìm tọa độ điểm $M$ thuộc trục $Oz$ sao cho $M$ cách đều hai điểm $A$ và $B$.}
  {$(0;0;3)$.}
  {$(0;3;0)$.}{$(-1;0;3)$.}{$(0;-2;3)$.}

\cauTN[id=C2-B8-D10-NB-019,muc=NB,loigiai={Gọi $M$ thuộc trục $Oz$, tọa độ $M$ chỉ khác $0$ ở thành phần tương ứng. Từ $MA=MB$ suy ra $MA^2=MB^2$, giải phương trình bậc nhất ta được $M=(0;0;-\dfrac{1}{4})$.}]{Trong không gian $Oxyz$, cho $A(1;-2;0)$, $B(-1;-1;-2)$. Tìm tọa độ điểm $M$ thuộc trục $Oz$ sao cho $M$ cách đều hai điểm $A$ và $B$.}
  {$(0;0;-\dfrac{1}{4})$.}
  {$(-2;0;-\dfrac{1}{4})$.}{$(2;0;-\dfrac{1}{4})$.}{$(0;0;\dfrac{3}{4})$.}

\cauTN[id=C2-B8-D10-NB-020,muc=NB,loigiai={Gọi $M$ thuộc trục $Oy$, tọa độ $M$ chỉ khác $0$ ở thành phần tương ứng. Từ $MA=MB$ suy ra $MA^2=MB^2$, giải phương trình bậc nhất ta được $M=(0;-\dfrac{4}{3};0)$.}]{Trong không gian $Oxyz$, cho $A(2;-3;0)$, $B(2;0;1)$. Tìm tọa độ điểm $M$ thuộc trục $Oy$ sao cho $M$ cách đều hai điểm $A$ và $B$.}
  {$(0;-\dfrac{4}{3};0)$.}
  {$(0;-\dfrac{10}{3};0)$.}{$(2;-\dfrac{4}{3};0)$.}{$(0;\dfrac{4}{3};0)$.}

\cauTN[id=C2-B8-D10-NB-021,muc=NB,loigiai={Gọi $M$ thuộc trục $Ox$, tọa độ $M$ chỉ khác $0$ ở thành phần tương ứng. Từ $MA=MB$ suy ra $MA^2=MB^2$, giải phương trình bậc nhất ta được $M=(0;0;0)$.}]{Trong không gian $Oxyz$, cho $A(3;0;-1)$, $B(-1;0;-3)$. Tìm tọa độ điểm $M$ thuộc trục $Ox$ sao cho $M$ cách đều hai điểm $A$ và $B$.}
  {$(0;0;0)$.}
  {$(-1;0;0)$.}{$(2;0;0)$.}{$(0;-2;0)$.}

\cauTN[id=C2-B8-D10-NB-022,muc=NB,loigiai={Gọi $M$ thuộc trục $Oy$, tọa độ $M$ chỉ khác $0$ ở thành phần tương ứng. Từ $MA=MB$ suy ra $MA^2=MB^2$, giải phương trình bậc nhất ta được $M=(0;\dfrac{7}{3};0)$.}]{Trong không gian $Oxyz$, cho $A(2;0;-2)$, $B(-2;3;3)$. Tìm tọa độ điểm $M$ thuộc trục $Oy$ sao cho $M$ cách đều hai điểm $A$ và $B$.}
  {$(0;\dfrac{7}{3};0)$.}
  {$(0;\dfrac{7}{3};-2)$.}{$(0;\dfrac{1}{3};0)$.}{$(\dfrac{7}{3};0;0)$.}

\cauTN[id=C2-B8-D10-NB-023,muc=NB,loigiai={Gọi $M$ thuộc trục $Ox$, tọa độ $M$ chỉ khác $0$ ở thành phần tương ứng. Từ $MA=MB$ suy ra $MA^2=MB^2$, giải phương trình bậc nhất ta được $M=(-\dfrac{9}{4};0;0)$.}]{Trong không gian $Oxyz$, cho $A(-3;3;0)$, $B(-1;2;-2)$. Tìm tọa độ điểm $M$ thuộc trục $Ox$ sao cho $M$ cách đều hai điểm $A$ và $B$.}
  {$(-\dfrac{9}{4};0;0)$.}
  {$(-\dfrac{13}{4};0;0)$.}{$(0;-\dfrac{9}{4};0)$.}{$(-\dfrac{9}{4};0;-1)$.}

\cauTN[id=C2-B8-D10-NB-024,muc=NB,loigiai={Gọi $M$ thuộc trục $Oz$, tọa độ $M$ chỉ khác $0$ ở thành phần tương ứng. Từ $MA=MB$ suy ra $MA^2=MB^2$, giải phương trình bậc nhất ta được $M=(0;0;1)$.}]{Trong không gian $Oxyz$, cho $A(-1;1;0)$, $B(1;-1;2)$. Tìm tọa độ điểm $M$ thuộc trục $Oz$ sao cho $M$ cách đều hai điểm $A$ và $B$.}
  {$(0;0;1)$.}
  {$(-2;0;1)$.}{$(0;2;1)$.}{$(1;0;1)$.}

\cauTN[id=C2-B8-D10-NB-025,muc=NB,loigiai={Với hệ trục đã chọn: $O(0;0;0)$, $A(4;0;0)$, $C(0;3;0)$, $O'(0;0;1)$ nên $B(4;3;0)$, $A'(4;0;1)$, $B'(4;3;1)$, $C'(0;3;1)$. Từ đó trọng tâm $G$ của tam giác $C'AA'$ có tọa độ $(\dfrac{8}{3};1;\dfrac{2}{3})$.}]{Cho hình hộp chữ nhật $OABC.O'A'B'C'$ có $OA=4$, $OC=3$, $OO'=1$. Chọn hệ trục tọa độ $Oxyz$ sao cho $A\in Ox$, $C\in Oy$, $O'\in Oz$. Tọa độ trọng tâm $G$ của tam giác $C'AA'$ là}
  {$(\dfrac{8}{3};1;\dfrac{2}{3})$.}
  {$(\dfrac{8}{3};1;-\dfrac{4}{3})$.}{$(\dfrac{14}{3};1;\dfrac{2}{3})$.}{$(\dfrac{11}{3};1;\dfrac{2}{3})$.}

\cauTN[id=C2-B8-D10-NB-026,muc=NB,loigiai={Với hệ trục đã chọn: $O(0;0;0)$, $A(3;0;0)$, $C(0;5;0)$, $O'(0;0;1)$ nên $B(3;5;0)$, $A'(3;0;1)$, $B'(3;5;1)$, $C'(0;5;1)$. Từ đó điểm $B$ có tọa độ $(3;5;0)$.}]{Cho hình hộp chữ nhật $OABC.O'A'B'C'$ có $OA=3$, $OC=5$, $OO'=1$. Chọn hệ trục tọa độ $Oxyz$ sao cho $A\in Ox$, $C\in Oy$, $O'\in Oz$. Tọa độ điểm $B$ là}
  {$(3;5;0)$.}
  {$(0;5;3)$.}{$(3;7;0)$.}{$(2;5;0)$.}

\cauTN[id=C2-B8-D10-NB-027,muc=NB,loigiai={Với hệ trục đã chọn: $O(0;0;0)$, $A(5;0;0)$, $C(0;2;0)$, $O'(0;0;2)$ nên $B(5;2;0)$, $A'(5;0;2)$, $B'(5;2;2)$, $C'(0;2;2)$. Từ đó trung điểm $I$ của đoạn thẳng $BO'$ có tọa độ $(\dfrac{5}{2};1;1)$.}]{Cho hình hộp chữ nhật $OABC.O'A'B'C'$ có $OA=5$, $OC=2$, $OO'=2$. Chọn hệ trục tọa độ $Oxyz$ sao cho $A\in Ox$, $C\in Oy$, $O'\in Oz$. Tọa độ trung điểm $I$ của đoạn thẳng $BO'$ là}
  {$(\dfrac{5}{2};1;1)$.}
  {$(\dfrac{3}{2};1;1)$.}{$(\dfrac{5}{2};-1;1)$.}{$(1;\dfrac{5}{2};1)$.}

\cauTN[id=C2-B8-D10-NB-028,muc=NB,loigiai={Với hệ trục đã chọn: $O(0;0;0)$, $A(4;0;0)$, $C(0;2;0)$, $O'(0;0;4)$ nên $B(4;2;0)$, $A'(4;0;4)$, $B'(4;2;4)$, $C'(0;2;4)$. Từ đó trọng tâm $G$ của tam giác $AOC'$ có tọa độ $(\dfrac{4}{3};\dfrac{2}{3};\dfrac{4}{3})$.}]{Cho hình hộp chữ nhật $OABC.O'A'B'C'$ có $OA=4$, $OC=2$, $OO'=4$. Chọn hệ trục tọa độ $Oxyz$ sao cho $A\in Ox$, $C\in Oy$, $O'\in Oz$. Tọa độ trọng tâm $G$ của tam giác $AOC'$ là}
  {$(\dfrac{4}{3};\dfrac{2}{3};\dfrac{4}{3})$.}
  {$(\dfrac{4}{3};\dfrac{5}{3};\dfrac{4}{3})$.}{$(-\dfrac{4}{3};\dfrac{2}{3};\dfrac{4}{3})$.}{$(\dfrac{4}{3};\dfrac{2}{3};-\dfrac{4}{3})$.}

\cauTN[id=C2-B8-D10-NB-029,muc=NB,loigiai={Với hệ trục đã chọn: $O(0;0;0)$, $A(2;0;0)$, $C(0;5;0)$, $O'(0;0;1)$ nên $B(2;5;0)$, $A'(2;0;1)$, $B'(2;5;1)$, $C'(0;5;1)$. Từ đó trọng tâm $G$ của tam giác $BCB'$ có tọa độ $(\dfrac{4}{3};5;\dfrac{1}{3})$.}]{Cho hình hộp chữ nhật $OABC.O'A'B'C'$ có $OA=2$, $OC=5$, $OO'=1$. Chọn hệ trục tọa độ $Oxyz$ sao cho $A\in Ox$, $C\in Oy$, $O'\in Oz$. Tọa độ trọng tâm $G$ của tam giác $BCB'$ là}
  {$(\dfrac{4}{3};5;\dfrac{1}{3})$.}
  {$(-\dfrac{2}{3};5;\dfrac{1}{3})$.}{$(5;\dfrac{4}{3};\dfrac{1}{3})$.}{$(\dfrac{1}{3};5;\dfrac{1}{3})$.}

\cauTN[id=C2-B8-D10-NB-030,muc=NB,loigiai={Với hệ trục đã chọn: $O(0;0;0)$, $A(3;0;0)$, $C(0;5;0)$, $O'(0;0;2)$ nên $B(3;5;0)$, $A'(3;0;2)$, $B'(3;5;2)$, $C'(0;5;2)$. Từ đó trung điểm $I$ của đoạn thẳng $C'A$ có tọa độ $(\dfrac{3}{2};\dfrac{5}{2};1)$.}]{Cho hình hộp chữ nhật $OABC.O'A'B'C'$ có $OA=3$, $OC=5$, $OO'=2$. Chọn hệ trục tọa độ $Oxyz$ sao cho $A\in Ox$, $C\in Oy$, $O'\in Oz$. Tọa độ trung điểm $I$ của đoạn thẳng $C'A$ là}
  {$(\dfrac{3}{2};\dfrac{5}{2};1)$.}
  {$(\dfrac{3}{2};\dfrac{5}{2};2)$.}{$(\dfrac{3}{2};\dfrac{7}{2};1)$.}{$(\dfrac{3}{2};\dfrac{3}{2};1)$.}

\cauTN[id=C2-B8-D10-NB-031,muc=NB,loigiai={Với hệ trục: $A(0;0;0)$, $B(5;0;0)$, $D(0;3;0)$, $S(0;0;3)$ nên $C(5;3;0)$. Do đó trung điểm $I$ của đoạn thẳng $SB$ có tọa độ $(\dfrac{5}{2};0;\dfrac{3}{2})$.}]{Cho hình chóp $S.ABCD$ có đáy $ABCD$ là hình chữ nhật, $AB=5$, $AD=3$, $SA\perp(ABCD)$ và $SA=3$. Chọn hệ trục tọa độ $Oxyz$ có gốc $A$, tia $AB$ là $Ox$, tia $AD$ là $Oy$, tia $AS$ là $Oz$. Tọa độ trung điểm $I$ của đoạn thẳng $SB$ là}
  {$(\dfrac{5}{2};0;\dfrac{3}{2})$.}
  {$(-\dfrac{5}{2};0;\dfrac{3}{2})$.}{$(0;\dfrac{5}{2};\dfrac{3}{2})$.}{$(\dfrac{3}{2};0;\dfrac{3}{2})$.}

\cauTN[id=C2-B8-D10-NB-032,muc=NB,loigiai={Với hệ trục: $A(0;0;0)$, $B(4;0;0)$, $D(0;3;0)$, $S(0;0;1)$ nên $C(4;3;0)$. Do đó trọng tâm $G$ của tam giác $DAS$ có tọa độ $(0;1;\dfrac{1}{3})$.}]{Cho hình chóp $S.ABCD$ có đáy $ABCD$ là hình chữ nhật, $AB=4$, $AD=3$, $SA\perp(ABCD)$ và $SA=1$. Chọn hệ trục tọa độ $Oxyz$ có gốc $A$, tia $AB$ là $Ox$, tia $AD$ là $Oy$, tia $AS$ là $Oz$. Tọa độ trọng tâm $G$ của tam giác $DAS$ là}
  {$(0;1;\dfrac{1}{3})$.}
  {$(0;1;\dfrac{7}{3})$.}{$(0;1;\dfrac{4}{3})$.}{$(-2;1;\dfrac{1}{3})$.}

\cauTN[id=C2-B8-D10-NB-033,muc=NB,loigiai={Với hệ trục: $A(0;0;0)$, $B(1;0;0)$, $D(0;2;0)$, $S(0;0;2)$ nên $C(1;2;0)$. Do đó điểm $C$ có tọa độ $(1;2;0)$.}]{Cho hình chóp $S.ABCD$ có đáy $ABCD$ là hình chữ nhật, $AB=1$, $AD=2$, $SA\perp(ABCD)$ và $SA=2$. Chọn hệ trục tọa độ $Oxyz$ có gốc $A$, tia $AB$ là $Ox$, tia $AD$ là $Oy$, tia $AS$ là $Oz$. Tọa độ điểm $C$ là}
  {$(1;2;0)$.}
  {$(1;0;0)$.}{$(1;2;-1)$.}{$(-1;2;0)$.}

\cauTN[id=C2-B8-D10-NB-034,muc=NB,loigiai={Với hệ trục: $A(0;0;0)$, $B(3;0;0)$, $D(0;2;0)$, $S(0;0;5)$ nên $C(3;2;0)$. Do đó trọng tâm $G$ của tam giác $ABS$ có tọa độ $(1;0;\dfrac{5}{3})$.}]{Cho hình chóp $S.ABCD$ có đáy $ABCD$ là hình chữ nhật, $AB=3$, $AD=2$, $SA\perp(ABCD)$ và $SA=5$. Chọn hệ trục tọa độ $Oxyz$ có gốc $A$, tia $AB$ là $Ox$, tia $AD$ là $Oy$, tia $AS$ là $Oz$. Tọa độ trọng tâm $G$ của tam giác $ABS$ là}
  {$(1;0;\dfrac{5}{3})$.}
  {$(-1;0;-\dfrac{5}{3})$.}{$(0;1;\dfrac{5}{3})$.}{$(1;0;-\dfrac{5}{3})$.}

\cauTN[id=C2-B8-D10-NB-035,muc=NB,loigiai={Với hệ trục: $A(0;0;0)$, $B(3;0;0)$, $D(0;2;0)$, $S(0;0;2)$ nên $C(3;2;0)$. Do đó trung điểm $I$ của đoạn thẳng $DA$ có tọa độ $(0;1;0)$.}]{Cho hình chóp $S.ABCD$ có đáy $ABCD$ là hình chữ nhật, $AB=3$, $AD=2$, $SA\perp(ABCD)$ và $SA=2$. Chọn hệ trục tọa độ $Oxyz$ có gốc $A$, tia $AB$ là $Ox$, tia $AD$ là $Oy$, tia $AS$ là $Oz$. Tọa độ trung điểm $I$ của đoạn thẳng $DA$ là}
  {$(0;1;0)$.}
  {$(0;0;0)$.}{$(2;1;0)$.}{$(1;1;0)$.}

\cauTN[id=C2-B8-D10-NB-036,muc=NB,loigiai={Với hệ trục: $A(0;0;0)$, $B(5;0;0)$, $D(0;2;0)$, $S(0;0;2)$ nên $C(5;2;0)$. Do đó trọng tâm $G$ của tam giác $ADB$ có tọa độ $(\dfrac{5}{3};\dfrac{2}{3};0)$.}]{Cho hình chóp $S.ABCD$ có đáy $ABCD$ là hình chữ nhật, $AB=5$, $AD=2$, $SA\perp(ABCD)$ và $SA=2$. Chọn hệ trục tọa độ $Oxyz$ có gốc $A$, tia $AB$ là $Ox$, tia $AD$ là $Oy$, tia $AS$ là $Oz$. Tọa độ trọng tâm $G$ của tam giác $ADB$ là}
  {$(\dfrac{5}{3};\dfrac{2}{3};0)$.}
  {$(-\dfrac{1}{3};\dfrac{2}{3};0)$.}{$(\dfrac{8}{3};\dfrac{2}{3};0)$.}{$(0;\dfrac{2}{3};\dfrac{5}{3})$.}

\cauDS[id=C2-B8-D10-NB-037,muc=NB,loigiai={Kiểm tra từng mệnh đề bằng công thức tọa độ. a) Đúng. b) Sai. c) Đúng. d) Đúng.}]{Trong không gian $Oxyz$, cho $A(3;2;-2)$, $B(3;-1;2)$, $C(1;3;0)$. Xét tính đúng sai của các mệnh đề sau:}{%
  \y{Độ dài đoạn thẳng $AB$ bằng $5$.}{d}
  \y{Diện tích tam giác $ABC$ bằng $\sqrt{51}$.}{s}
  \y{Điểm $M$ thuộc trục $Oy$ cách đều $A$ và $B$ có tọa độ $(0;\dfrac{1}{2};0)$.}{d}
  \y{Diện tích tam giác $ABC$ bằng $5\sqrt{2}$.}{d}}

\cauDS[id=C2-B8-D10-NB-038,muc=NB,loigiai={Kiểm tra từng mệnh đề bằng công thức tọa độ. a) Đúng. b) Sai. c) Đúng. d) Đúng.}]{Trong không gian $Oxyz$, cho $A(2;-1;0)$, $B(-1;-3;2)$, $C(-1;-3;-3)$. Xét tính đúng sai của các mệnh đề sau:}{%
  \y{Diện tích tam giác $ABC$ bằng $\dfrac{5\sqrt{13}}{2}$.}{d}
  \y{Trọng tâm tam giác $ABC$ có tọa độ $(0;-\dfrac{7}{3};\dfrac{5}{3})$.}{s}
  \y{Trọng tâm tam giác $ABC$ có tọa độ $(0;-\dfrac{7}{3};-\dfrac{1}{3})$.}{d}
  \y{Độ dài đoạn thẳng $AB$ bằng $\sqrt{17}$.}{d}}

\cauDS[id=C2-B8-D10-NB-039,muc=NB,loigiai={Kiểm tra từng mệnh đề bằng công thức tọa độ. a) Đúng. b) Sai. c) Đúng. d) Sai.}]{Trong không gian $Oxyz$, cho $A(1;3;1)$, $B(-2;-1;2)$, $C(-2;-3;0)$. Xét tính đúng sai của các mệnh đề sau:}{%
  \y{Trọng tâm tam giác $ABC$ có tọa độ $(-1;-\dfrac{1}{3};1)$.}{d}
  \y{Trọng tâm tam giác $ABC$ có tọa độ $(-1;-\dfrac{4}{3};1)$.}{s}
  \y{Độ dài đoạn thẳng $AB$ bằng $\sqrt{26}$.}{d}
  \y{Diện tích tam giác $ABC$ bằng $2\sqrt{11}$.}{s}}

\cauDS[id=C2-B8-D10-NB-040,muc=NB,loigiai={Kiểm tra từng mệnh đề bằng công thức tọa độ. a) Sai. b) Sai. c) Đúng. d) Sai.}]{Trong không gian $Oxyz$, cho $A(-3;1;0)$, $B(3;-2;0)$, $C(3;3;-1)$. Xét tính đúng sai của các mệnh đề sau:}{%
  \y{Trọng tâm tam giác $ABC$ có tọa độ $(1;-\dfrac{1}{3};-\dfrac{1}{3})$.}{s}
  \y{Điểm $M$ thuộc trục $Ox$ cách đều $A$ và $B$ có tọa độ $(0;\dfrac{1}{4};0)$.}{s}
  \y{Diện tích tam giác $ABC$ bằng $\dfrac{3\sqrt{105}}{2}$.}{d}
  \y{Độ dài đoạn thẳng $AB$ bằng $2\sqrt{11}$.}{s}}

\cauDS[id=C2-B8-D10-NB-041,muc=NB,loigiai={Kiểm tra từng mệnh đề bằng công thức tọa độ. a) Đúng. b) Sai. c) Đúng. d) Đúng.}]{Trong không gian $Oxyz$, cho $A(-1;0;-1)$, $B(-2;-2;-3)$, $C(1;0;0)$. Xét tính đúng sai của các mệnh đề sau:}{%
  \y{Điểm $M$ thuộc trục $Oz$ cách đều $A$ và $B$ có tọa độ $(0;0;-\dfrac{15}{4})$.}{d}
  \y{Diện tích tam giác $ABC$ bằng $\sqrt{29}$.}{s}
  \y{Diện tích tam giác $ABC$ bằng $\dfrac{\sqrt{29}}{2}$.}{d}
  \y{Trọng tâm tam giác $ABC$ có tọa độ $(-\dfrac{2}{3};-\dfrac{2}{3};-\dfrac{4}{3})$.}{d}}

\cauDS[id=C2-B8-D10-NB-042,muc=NB,loigiai={Kiểm tra từng mệnh đề bằng công thức tọa độ. a) Sai. b) Đúng. c) Sai. d) Sai.}]{Trong không gian $Oxyz$, cho $A(2;1;-1)$, $B(-1;-1;3)$, $C(3;-2;-1)$. Xét tính đúng sai của các mệnh đề sau:}{%
  \y{Trọng tâm tam giác $ABC$ có tọa độ $(\dfrac{4}{3};-\dfrac{2}{3};-\dfrac{1}{3})$.}{s}
  \y{Diện tích tam giác $ABC$ bằng $\dfrac{\sqrt{281}}{2}$.}{d}
  \y{Điểm $M$ thuộc trục $Oy$ cách đều $A$ và $B$ có tọa độ $(2;-\dfrac{5}{4};0)$.}{s}
  \y{Điểm $M$ thuộc trục $Oy$ cách đều $A$ và $B$ có tọa độ $(0;-\dfrac{5}{4};2)$.}{s}}

\cauDS[id=C2-B8-D10-NB-043,muc=NB,loigiai={Kiểm tra từng mệnh đề bằng công thức tọa độ. a) Sai. b) Đúng. c) Sai. d) Sai.}]{Trong không gian $Oxyz$, cho $A(3;-1;1)$, $B(2;-2;-2)$, $C(2;1;0)$. Xét tính đúng sai của các mệnh đề sau:}{%
  \y{Điểm $M$ thuộc trục $Oy$ cách đều $A$ và $B$ có tọa độ $(0;-\dfrac{3}{2};0)$.}{s}
  \y{Điểm $M$ thuộc trục $Oy$ cách đều $A$ và $B$ có tọa độ $(0;-\dfrac{1}{2};0)$.}{d}
  \y{Diện tích tam giác $ABC$ bằng $\sqrt{62}$.}{s}
  \y{Độ dài đoạn thẳng $AB$ bằng $\sqrt{10}$.}{s}}

\cauDS[id=C2-B8-D10-NB-044,muc=NB,loigiai={Kiểm tra từng mệnh đề bằng công thức tọa độ. a) Sai. b) Đúng. c) Sai. d) Sai.}]{Trong không gian $Oxyz$, cho $A(-3;0;3)$, $B(3;-1;-3)$, $C(2;0;0)$. Xét tính đúng sai của các mệnh đề sau:}{%
  \y{Trọng tâm tam giác $ABC$ có tọa độ $(\dfrac{2}{3};0;-\dfrac{1}{3})$.}{s}
  \y{Diện tích tam giác $ABC$ bằng $\dfrac{\sqrt{178}}{2}$.}{d}
  \y{Độ dài đoạn thẳng $AB$ bằng $\sqrt{77}$.}{s}
  \y{Diện tích tam giác $ABC$ bằng $\dfrac{\sqrt{182}}{2}$.}{s}}

\cauDS[id=C2-B8-D10-NB-045,muc=NB,loigiai={Kiểm tra từng mệnh đề bằng công thức tọa độ. a) Sai. b) Sai. c) Sai. d) Đúng.}]{Trong không gian $Oxyz$, cho $A(-2;1;-3)$, $B(-1;-1;-1)$, $C(1;-2;-2)$. Xét tính đúng sai của các mệnh đề sau:}{%
  \y{Diện tích tam giác $ABC$ bằng $5\sqrt{2}$.}{s}
  \y{Độ dài đoạn thẳng $AB$ bằng $2\sqrt{2}$.}{s}
  \y{Trọng tâm tam giác $ABC$ có tọa độ $(\dfrac{4}{3};-\dfrac{2}{3};-2)$.}{s}
  \y{Điểm $M$ thuộc trục $Oz$ cách đều $A$ và $B$ có tọa độ $(0;0;-\dfrac{11}{4})$.}{d}}

\cauDS[id=C2-B8-D10-NB-046,muc=NB,loigiai={Kiểm tra từng mệnh đề bằng công thức tọa độ. a) Đúng. b) Đúng. c) Sai. d) Sai.}]{Trong không gian $Oxyz$, cho $A(1;0;2)$, $B(-3;2;1)$, $C(-3;2;0)$. Xét tính đúng sai của các mệnh đề sau:}{%
  \y{Độ dài đoạn thẳng $AB$ bằng $\sqrt{21}$.}{d}
  \y{Trọng tâm tam giác $ABC$ có tọa độ $(-\dfrac{5}{3};\dfrac{4}{3};1)$.}{d}
  \y{Diện tích tam giác $ABC$ bằng $\sqrt{10}$.}{s}
  \y{Trọng tâm tam giác $ABC$ có tọa độ $(-\dfrac{5}{3};\dfrac{4}{3};3)$.}{s}}

\cauDS[id=C2-B8-D10-NB-047,muc=NB,loigiai={Kiểm tra từng mệnh đề bằng công thức tọa độ. a) Sai. b) Đúng. c) Đúng. d) Sai.}]{Trong không gian $Oxyz$, cho $A(3;3;-1)$, $B(-1;3;0)$, $C(2;-1;1)$. Xét tính đúng sai của các mệnh đề sau:}{%
  \y{Điểm $M$ thuộc trục $Ox$ cách đều $A$ và $B$ có tọa độ $(\dfrac{9}{8};-2;0)$.}{s}
  \y{Điểm $M$ thuộc trục $Ox$ cách đều $A$ và $B$ có tọa độ $(\dfrac{9}{8};0;0)$.}{d}
  \y{Diện tích tam giác $ABC$ bằng $\dfrac{\sqrt{321}}{2}$.}{d}
  \y{Trọng tâm tam giác $ABC$ có tọa độ $(\dfrac{1}{3};\dfrac{5}{3};0)$.}{s}}

\cauDS[id=C2-B8-D10-NB-048,muc=NB,loigiai={Kiểm tra từng mệnh đề bằng công thức tọa độ. a) Đúng. b) Sai. c) Sai. d) Sai.}]{Trong không gian $Oxyz$, cho $A(-2;3;0)$, $B(2;-2;-1)$, $C(3;2;1)$. Xét tính đúng sai của các mệnh đề sau:}{%
  \y{Độ dài đoạn thẳng $AB$ bằng $\sqrt{42}$.}{d}
  \y{Trọng tâm tam giác $ABC$ có tọa độ $(\dfrac{3}{2};\dfrac{3}{2};0)$.}{s}
  \y{Điểm $M$ thuộc trục $Oz$ cách đều $A$ và $B$ có tọa độ $(0;0;3)$.}{s}
  \y{Trọng tâm tam giác $ABC$ có tọa độ $(-1;1;0)$.}{s}}

\cauTLN[id=C2-B8-D10-NB-049,muc=NB,loigiai={Ta có $\overrightarrow{AB}=(-1;2;2)$, $\overrightarrow{AC}=(3;0;1)$. Diện tích $S=\dfrac{1}{2}\sqrt{|\overrightarrow{AB}|^2|\overrightarrow{AC}|^2-(\overrightarrow{AB}\cdot\overrightarrow{AC})^2}=\dfrac{1}{2}\sqrt{9\cdot10+1^2}=\dfrac{\sqrt{89}}{2}$.}]{Trong không gian $Oxyz$, cho $A(-2;-1;-2)$, $B(-3;1;0)$, $C(1;-1;-1)$. Tính $4S^2$, trong đó $S$ là diện tích tam giác $ABC$.}{89}

\cauTLN[id=C2-B8-D10-NB-050,muc=NB,loigiai={Ta có $\overrightarrow{AB}=(-2;4;0)$, $\overrightarrow{AC}=(-1;0;-2)$. Diện tích $S=\dfrac{1}{2}\sqrt{|\overrightarrow{AB}|^2|\overrightarrow{AC}|^2-(\overrightarrow{AB}\cdot\overrightarrow{AC})^2}=\dfrac{1}{2}\sqrt{20\cdot5-2^2}=2\sqrt{6}$.}]{Trong không gian $Oxyz$, cho $A(0;-3;0)$, $B(-2;1;0)$, $C(-1;-3;-2)$. Tính $4S^2$, trong đó $S$ là diện tích tam giác $ABC$.}{96}

\cauTLN[id=C2-B8-D10-NB-051,muc=NB,loigiai={Ta có $\overrightarrow{AB}=(-5;0;-5)$, $\overrightarrow{AC}=(-2;-1;-2)$. Diện tích $S=\dfrac{1}{2}\sqrt{|\overrightarrow{AB}|^2|\overrightarrow{AC}|^2-(\overrightarrow{AB}\cdot\overrightarrow{AC})^2}=\dfrac{1}{2}\sqrt{50\cdot9-20^2}=\dfrac{5\sqrt{2}}{2}$.}]{Trong không gian $Oxyz$, cho $A(3;3;2)$, $B(-2;3;-3)$, $C(1;2;0)$. Tính $4S^2$, trong đó $S$ là diện tích tam giác $ABC$.}{50}

\cauTLN[id=C2-B8-D10-NB-052,muc=NB,loigiai={Ta có $\overrightarrow{AB}=(1;-3;1)$, $\overrightarrow{AC}=(-2;0;1)$. Diện tích $S=\dfrac{1}{2}\sqrt{|\overrightarrow{AB}|^2|\overrightarrow{AC}|^2-(\overrightarrow{AB}\cdot\overrightarrow{AC})^2}=\dfrac{1}{2}\sqrt{11\cdot5+1^2}=\dfrac{3\sqrt{6}}{2}$.}]{Trong không gian $Oxyz$, cho $A(-1;0;-1)$, $B(0;-3;0)$, $C(-3;0;0)$. Tính $4S^2$, trong đó $S$ là diện tích tam giác $ABC$.}{54}

\cauTLN[id=C2-B8-D10-NB-053,muc=NB,loigiai={Ta có $\overrightarrow{AB}=(2;-1;1)$, $\overrightarrow{AC}=(1;0;0)$. Diện tích $S=\dfrac{1}{2}\sqrt{|\overrightarrow{AB}|^2|\overrightarrow{AC}|^2-(\overrightarrow{AB}\cdot\overrightarrow{AC})^2}=\dfrac{1}{2}\sqrt{6\cdot1-2^2}=\dfrac{\sqrt{2}}{2}$.}]{Trong không gian $Oxyz$, cho $A(-1;1;-3)$, $B(1;0;-2)$, $C(0;1;-3)$. Tính $4S^2$, trong đó $S$ là diện tích tam giác $ABC$.}{2}

\cauTLN[id=C2-B8-D10-NB-054,muc=NB,loigiai={Ta có $\overrightarrow{AB}=(-1;-4;2)$, $\overrightarrow{AC}=(0;0;3)$. Diện tích $S=\dfrac{1}{2}\sqrt{|\overrightarrow{AB}|^2|\overrightarrow{AC}|^2-(\overrightarrow{AB}\cdot\overrightarrow{AC})^2}=\dfrac{1}{2}\sqrt{21\cdot9-6^2}=\dfrac{3\sqrt{17}}{2}$.}]{Trong không gian $Oxyz$, cho $A(2;2;-3)$, $B(1;-2;-1)$, $C(2;2;0)$. Tính $4S^2$, trong đó $S$ là diện tích tam giác $ABC$.}{153}

\cauTLN[id=C2-B8-D10-NB-055,muc=NB,loigiai={Gọi $M$ thuộc trục $Oz$, tọa độ $M$ chỉ khác $0$ ở thành phần tương ứng. Từ $MA=MB$ suy ra $MA^2=MB^2$, giải phương trình bậc nhất ta được $M=(0;0;0)$.}]{Trong không gian $Oxyz$, cho $A(0;-2;-3)$, $B(3;0;-2)$. Điểm $M$ thuộc trục $Oz$ cách đều hai điểm $A$ và $B$. Tính tọa độ khác $0$ của $M$ nhân với $2$ (kết quả nguyên).}{0}

\cauTLN[id=C2-B8-D10-NB-056,muc=NB,loigiai={Gọi $M$ thuộc trục $Oz$, tọa độ $M$ chỉ khác $0$ ở thành phần tương ứng. Từ $MA=MB$ suy ra $MA^2=MB^2$, giải phương trình bậc nhất ta được $M=(0;0;-1)$.}]{Trong không gian $Oxyz$, cho $A(-2;-1;-1)$, $B(0;2;-2)$. Điểm $M$ thuộc trục $Oz$ cách đều hai điểm $A$ và $B$. Tính tọa độ khác $0$ của $M$ nhân với $2$ (kết quả nguyên).}{-2}

\cauTLN[id=C2-B8-D10-NB-057,muc=NB,loigiai={Gọi $M$ thuộc trục $Ox$, tọa độ $M$ chỉ khác $0$ ở thành phần tương ứng. Từ $MA=MB$ suy ra $MA^2=MB^2$, giải phương trình bậc nhất ta được $M=(2;0;0)$.}]{Trong không gian $Oxyz$, cho $A(1;-3;-3)$, $B(3;3;3)$. Điểm $M$ thuộc trục $Ox$ cách đều hai điểm $A$ và $B$. Tính tọa độ khác $0$ của $M$ nhân với $2$ (kết quả nguyên).}{4}

\cauTLN[id=C2-B8-D10-NB-058,muc=NB,loigiai={Gọi $M$ thuộc trục $Oy$, tọa độ $M$ chỉ khác $0$ ở thành phần tương ứng. Từ $MA=MB$ suy ra $MA^2=MB^2$, giải phương trình bậc nhất ta được $M=(0;0;0)$.}]{Trong không gian $Oxyz$, cho $A(-2;-1;-2)$, $B(0;3;0)$. Điểm $M$ thuộc trục $Oy$ cách đều hai điểm $A$ và $B$. Tính tọa độ khác $0$ của $M$ nhân với $2$ (kết quả nguyên).}{0}

\cauTLN[id=C2-B8-D10-NB-059,muc=NB,loigiai={Gọi $M$ thuộc trục $Ox$, tọa độ $M$ chỉ khác $0$ ở thành phần tương ứng. Từ $MA=MB$ suy ra $MA^2=MB^2$, giải phương trình bậc nhất ta được $M=(2;0;0)$.}]{Trong không gian $Oxyz$, cho $A(0;0;-2)$, $B(2;2;2)$. Điểm $M$ thuộc trục $Ox$ cách đều hai điểm $A$ và $B$. Tính tọa độ khác $0$ của $M$ nhân với $2$ (kết quả nguyên).}{4}

\cauTLN[id=C2-B8-D10-NB-060,muc=NB,loigiai={Gọi $M$ thuộc trục $Oz$, tọa độ $M$ chỉ khác $0$ ở thành phần tương ứng. Từ $MA=MB$ suy ra $MA^2=MB^2$, giải phương trình bậc nhất ta được $M=(0;0;3)$.}]{Trong không gian $Oxyz$, cho $A(1;3;3)$, $B(0;-1;0)$. Điểm $M$ thuộc trục $Oz$ cách đều hai điểm $A$ và $B$. Tính tọa độ khác $0$ của $M$ nhân với $2$ (kết quả nguyên).}{6}

\cauTN[id=C2-B8-D10-TH-001,muc=TH,loigiai={$AB=1$, $AC=5$, $BC=4$ nên chu vi bằng $1+5+4=10$.}]{Tam giác $ABC$ có $A(3;-1;-1)$, $B(3;-1;0)$, $C(3;-1;4)$. Chu vi tam giác $ABC$ bằng}
  {$10$.}
  {$11$.}{$9$.}{$12$.}

\cauTN[id=C2-B8-D10-TH-002,muc=TH,loigiai={$AB=7$, $AC=3$, $BC=6$ nên chu vi bằng $7+3+6=16$.}]{Tam giác $ABC$ có $A(0;3;-1)$, $B(2;9;2)$, $C(-2;5;0)$. Chu vi tam giác $ABC$ bằng}
  {$16$.}
  {$17$.}{$15$.}{$18$.}

\cauTN[id=C2-B8-D10-TH-003,muc=TH,loigiai={$AB=3$, $AC=3$, $BC=2$ nên chu vi bằng $3+3+2=8$.}]{Tam giác $ABC$ có $A(3;0;2)$, $B(5;2;1)$, $C(5;2;3)$. Chu vi tam giác $ABC$ bằng}
  {$8$.}
  {$9$.}{$7$.}{$10$.}

\cauTN[id=C2-B8-D10-TH-004,muc=TH,loigiai={$AB=5$, $AC=1$, $BC=4$ nên chu vi bằng $5+1+4=10$.}]{Tam giác $ABC$ có $A(1;-2;-3)$, $B(1;-2;2)$, $C(1;-2;-2)$. Chu vi tam giác $ABC$ bằng}
  {$10$.}
  {$11$.}{$9$.}{$12$.}

\cauTN[id=C2-B8-D10-TH-005,muc=TH,loigiai={$AB=3$, $AC=3$, $BC=2$ nên chu vi bằng $3+3+2=8$.}]{Tam giác $ABC$ có $A(2;-3;-3)$, $B(4;-1;-2)$, $C(4;-1;-4)$. Chu vi tam giác $ABC$ bằng}
  {$8$.}
  {$9$.}{$7$.}{$10$.}

\cauTN[id=C2-B8-D10-TH-006,muc=TH,loigiai={$AB=5$, $AC=1$, $BC=4$ nên chu vi bằng $5+1+4=10$.}]{Tam giác $ABC$ có $A(-1;2;-3)$, $B(4;2;-3)$, $C(0;2;-3)$. Chu vi tam giác $ABC$ bằng}
  {$10$.}
  {$11$.}{$9$.}{$12$.}

\cauTN[id=C2-B8-D10-TH-007,muc=TH,loigiai={$AB=3$, $AC=7$, $BC=6$ nên chu vi bằng $3+7+6=16$.}]{Tam giác $ABC$ có $A(0;2;-1)$, $B(1;0;1)$, $C(3;4;5)$. Chu vi tam giác $ABC$ bằng}
  {$16$.}
  {$17$.}{$15$.}{$18$.}

\cauTN[id=C2-B8-D10-TH-008,muc=TH,loigiai={$AB=7$, $AC=3$, $BC=6$ nên chu vi bằng $7+3+6=16$.}]{Tam giác $ABC$ có $A(2;1;2)$, $B(5;3;8)$, $C(3;-1;4)$. Chu vi tam giác $ABC$ bằng}
  {$16$.}
  {$17$.}{$15$.}{$18$.}

\cauTN[id=C2-B8-D10-TH-009,muc=TH,loigiai={$AB=9$, $AC=5$, $BC=6$ nên chu vi bằng $9+5+6=20$.}]{Tam giác $ABC$ có $A(0;-3;-1)$, $B(4;1;6)$, $C(0;-3;4)$. Chu vi tam giác $ABC$ bằng}
  {$20$.}
  {$21$.}{$19$.}{$22$.}

\cauTN[id=C2-B8-D10-TH-010,muc=TH,loigiai={$AB=3$, $AC=3$, $BC=2$ nên chu vi bằng $3+3+2=8$.}]{Tam giác $ABC$ có $A(-3;-3;1)$, $B(-1;-4;3)$, $C(-1;-2;3)$. Chu vi tam giác $ABC$ bằng}
  {$8$.}
  {$9$.}{$7$.}{$10$.}

\cauTN[id=C2-B8-D10-TH-011,muc=TH,loigiai={$AB=3$, $AC=3$, $BC=2$ nên chu vi bằng $3+3+2=8$.}]{Tam giác $ABC$ có $A(1;0;-2)$, $B(3;2;-3)$, $C(3;2;-1)$. Chu vi tam giác $ABC$ bằng}
  {$8$.}
  {$9$.}{$7$.}{$10$.}

\cauTN[id=C2-B8-D10-TH-012,muc=TH,loigiai={$AB=5$, $AC=9$, $BC=6$ nên chu vi bằng $5+9+6=20$.}]{Tam giác $ABC$ có $A(0;-3;0)$, $B(0;-3;5)$, $C(4;1;7)$. Chu vi tam giác $ABC$ bằng}
  {$20$.}
  {$21$.}{$19$.}{$22$.}

\cauTN[id=C2-B8-D10-TH-013,muc=TH,loigiai={Tính độ dài các cạnh: $AB^2=8$, $AC^2=32$, $BC^2=40$ và các tích vô hướng. Chỉ có tam giác vuông tại $A$ là đúng.}]{Cho tam giác $ABC$ có $A(-3;3;1)$, $B(-1;3;3)$, $C(1;3;-3)$. Khẳng định nào dưới đây đúng?}
  {Tam giác $ABC$ vuông tại $A$.}
  {Tam giác $ABC$ cân tại $A$.}{Tam giác $ABC$ cân tại $B$.}{Tam giác $ABC$ cân tại $C$.}

\cauTN[id=C2-B8-D10-TH-014,muc=TH,loigiai={Tính độ dài các cạnh: $AB^2=9$, $AC^2=81$, $BC^2=72$ và các tích vô hướng. Chỉ có tam giác vuông tại $B$ là đúng.}]{Cho tam giác $ABC$ có $A(3;2;3)$, $B(3;-1;3)$, $C(-3;-1;-3)$. Khẳng định nào dưới đây đúng?}
  {Tam giác $ABC$ vuông tại $B$.}
  {Tam giác $ABC$ cân tại $B$.}{Tam giác $ABC$ vuông tại $C$.}{Tam giác $ABC$ vuông tại $A$.}

\cauTN[id=C2-B8-D10-TH-015,muc=TH,loigiai={Tính độ dài các cạnh: $AB^2=41$, $AC^2=51$, $BC^2=10$ và các tích vô hướng. Chỉ có tam giác vuông tại $B$ là đúng.}]{Cho tam giác $ABC$ có $A(-1;-3;2)$, $B(-2;3;0)$, $C(-2;2;-3)$. Khẳng định nào dưới đây đúng?}
  {Tam giác $ABC$ vuông tại $B$.}
  {Tam giác $ABC$ cân tại $A$.}{Tam giác $ABC$ vuông tại $C$.}{Tam giác $ABC$ cân tại $C$.}

\cauTN[id=C2-B8-D10-TH-016,muc=TH,loigiai={Tính độ dài các cạnh: $AB^2=6$, $AC^2=14$, $BC^2=20$ và các tích vô hướng. Chỉ có tam giác vuông tại $A$ là đúng.}]{Cho tam giác $ABC$ có $A(1;1;2)$, $B(2;-1;1)$, $C(2;3;-1)$. Khẳng định nào dưới đây đúng?}
  {Tam giác $ABC$ vuông tại $A$.}
  {Tam giác $ABC$ cân tại $B$.}{Tam giác $ABC$ vuông tại $B$.}{Tam giác $ABC$ vuông tại $C$.}

\cauTN[id=C2-B8-D10-TH-017,muc=TH,loigiai={Tính độ dài các cạnh: $AB^2=26$, $AC^2=8$, $BC^2=26$ và các tích vô hướng. Chỉ có tam giác cân tại $B$ là đúng.}]{Cho tam giác $ABC$ có $A(-2;0;-3)$, $B(-3;3;1)$, $C(0;2;-3)$. Khẳng định nào dưới đây đúng?}
  {Tam giác $ABC$ cân tại $B$.}
  {Tam giác $ABC$ vuông tại $C$.}{Tam giác $ABC$ vuông tại $B$.}{Tam giác $ABC$ cân tại $C$.}

\cauTN[id=C2-B8-D10-TH-018,muc=TH,loigiai={Tính độ dài các cạnh: $AB^2=6$, $AC^2=35$, $BC^2=29$ và các tích vô hướng. Chỉ có tam giác vuông tại $B$ là đúng.}]{Cho tam giác $ABC$ có $A(3;-2;-2)$, $B(2;0;-1)$, $C(-2;-3;1)$. Khẳng định nào dưới đây đúng?}
  {Tam giác $ABC$ vuông tại $B$.}
  {Tam giác $ABC$ cân tại $A$.}{Tam giác $ABC$ cân tại $B$.}{Tam giác $ABC$ cân tại $C$.}

\cauTN[id=C2-B8-D10-TH-019,muc=TH,loigiai={Tính độ dài các cạnh: $AB^2=19$, $AC^2=17$, $BC^2=2$ và các tích vô hướng. Chỉ có tam giác vuông tại $C$ là đúng.}]{Cho tam giác $ABC$ có $A(1;3;3)$, $B(2;0;0)$, $C(3;1;0)$. Khẳng định nào dưới đây đúng?}
  {Tam giác $ABC$ vuông tại $C$.}
  {Tam giác $ABC$ cân tại $B$.}{Tam giác $ABC$ vuông tại $A$.}{Tam giác $ABC$ cân tại $A$.}

\cauTN[id=C2-B8-D10-TH-020,muc=TH,loigiai={Tính độ dài các cạnh: $AB^2=32$, $AC^2=9$, $BC^2=41$ và các tích vô hướng. Chỉ có tam giác vuông tại $A$ là đúng.}]{Cho tam giác $ABC$ có $A(1;2;-3)$, $B(1;-2;1)$, $C(-2;2;-3)$. Khẳng định nào dưới đây đúng?}
  {Tam giác $ABC$ vuông tại $A$.}
  {Tam giác $ABC$ cân tại $C$.}{Tam giác $ABC$ vuông tại $C$.}{Tam giác $ABC$ cân tại $A$.}

\cauTN[id=C2-B8-D10-TH-021,muc=TH,loigiai={Tính độ dài các cạnh: $AB^2=27$, $AC^2=24$, $BC^2=27$ và các tích vô hướng. Chỉ có tam giác cân tại $B$ là đúng.}]{Cho tam giác $ABC$ có $A(-2;0;0)$, $B(1;3;-3)$, $C(2;-2;-2)$. Khẳng định nào dưới đây đúng?}
  {Tam giác $ABC$ cân tại $B$.}
  {Tam giác $ABC$ vuông tại $C$.}{Tam giác $ABC$ vuông tại $B$.}{Tam giác $ABC$ cân tại $C$.}

\cauTN[id=C2-B8-D10-TH-022,muc=TH,loigiai={Tính độ dài các cạnh: $AB^2=41$, $AC^2=8$, $BC^2=33$ và các tích vô hướng. Chỉ có tam giác vuông tại $C$ là đúng.}]{Cho tam giác $ABC$ có $A(-2;3;2)$, $B(2;3;-3)$, $C(0;1;2)$. Khẳng định nào dưới đây đúng?}
  {Tam giác $ABC$ vuông tại $C$.}
  {Tam giác $ABC$ cân tại $A$.}{Tam giác $ABC$ cân tại $B$.}{Tam giác $ABC$ vuông tại $A$.}

\cauTN[id=C2-B8-D10-TH-023,muc=TH,loigiai={Tính độ dài các cạnh: $AB^2=9$, $AC^2=9$, $BC^2=26$ và các tích vô hướng. Chỉ có tam giác cân tại $A$ là đúng.}]{Cho tam giác $ABC$ có $A(0;-2;0)$, $B(-2;-3;2)$, $C(-1;0;-2)$. Khẳng định nào dưới đây đúng?}
  {Tam giác $ABC$ cân tại $A$.}
  {Tam giác $ABC$ vuông tại $B$.}{Tam giác $ABC$ vuông tại $C$.}{Tam giác $ABC$ cân tại $C$.}

\cauTN[id=C2-B8-D10-TH-024,muc=TH,loigiai={Tính độ dài các cạnh: $AB^2=18$, $AC^2=17$, $BC^2=35$ và các tích vô hướng. Chỉ có tam giác vuông tại $A$ là đúng.}]{Cho tam giác $ABC$ có $A(-3;0;1)$, $B(0;3;1)$, $C(-1;-2;-2)$. Khẳng định nào dưới đây đúng?}
  {Tam giác $ABC$ vuông tại $A$.}
  {Tam giác $ABC$ vuông tại $C$.}{Tam giác $ABC$ cân tại $C$.}{Tam giác $ABC$ cân tại $B$.}

\cauTN[id=C2-B8-D10-TH-025,muc=TH,loigiai={Diện tích $S=\dfrac{1}{2}\sqrt{AB^2\cdot AC^2-(\overrightarrow{AB}\cdot\overrightarrow{AC})^2}$, $BC=\sqrt{77}$, đường cao $AH=\dfrac{2S}{BC}=\dfrac{\sqrt{4774}}{11}$.}]{Cho tam giác $ABC$ có $A(-1;-3;-4)$, $B(4;1;-3)$, $C(2;-2;5)$. Độ dài đường cao kẻ từ $A$ của tam giác bằng}
  {$\dfrac{\sqrt{4774}}{11}$.}
  {$\dfrac{\sqrt{4895}}{11}$.}{$\dfrac{2\sqrt{2387}}{11}$.}{$\dfrac{\sqrt{2387}}{11}$.}

\cauTN[id=C2-B8-D10-TH-026,muc=TH,loigiai={Diện tích $S=\dfrac{1}{2}\sqrt{AB^2\cdot AC^2-(\overrightarrow{AB}\cdot\overrightarrow{AC})^2}$, $BC=\sqrt{51}$, đường cao $AH=\dfrac{2S}{BC}=\dfrac{2\sqrt{25194}}{51}$.}]{Cho tam giác $ABC$ có $A(5;-3;2)$, $B(-2;0;-5)$, $C(-1;-1;2)$. Độ dài đường cao kẻ từ $A$ của tam giác bằng}
  {$\dfrac{2\sqrt{25194}}{51}$.}
  {$\dfrac{\sqrt{103377}}{51}$.}{$\dfrac{4\sqrt{12597}}{51}$.}{$\dfrac{2\sqrt{12597}}{51}$.}

\cauTN[id=C2-B8-D10-TH-027,muc=TH,loigiai={Diện tích $S=\dfrac{1}{2}\sqrt{AB^2\cdot AC^2-(\overrightarrow{AB}\cdot\overrightarrow{AC})^2}$, $BC=\sqrt{26}$, đường cao $AH=\dfrac{2S}{BC}=\dfrac{\sqrt{12233}}{13}$.}]{Cho tam giác $ABC$ có $A(-4;5;3)$, $B(-1;-3;3)$, $C(-5;-4;0)$. Độ dài đường cao kẻ từ $A$ của tam giác bằng}
  {$\dfrac{\sqrt{12233}}{13}$.}
  {$\dfrac{3\sqrt{1378}}{13}$.}{$\dfrac{\sqrt{24466}}{13}$.}{$\dfrac{\sqrt{24466}}{26}$.}

\cauTN[id=C2-B8-D10-TH-028,muc=TH,loigiai={Diện tích $S=\dfrac{1}{2}\sqrt{AB^2\cdot AC^2-(\overrightarrow{AB}\cdot\overrightarrow{AC})^2}$, $BC=\sqrt{13}$, đường cao $AH=\dfrac{2S}{BC}=\dfrac{6\sqrt{182}}{13}$.}]{Cho tam giác $ABC$ có $A(4;5;-4)$, $B(2;-1;-4)$, $C(0;-1;-1)$. Độ dài đường cao kẻ từ $A$ của tam giác bằng}
  {$\dfrac{6\sqrt{182}}{13}$.}
  {$\dfrac{\sqrt{6721}}{13}$.}{$\dfrac{12\sqrt{91}}{13}$.}{$\dfrac{6\sqrt{91}}{13}$.}

\cauTN[id=C2-B8-D10-TH-029,muc=TH,loigiai={Diện tích $S=\dfrac{1}{2}\sqrt{AB^2\cdot AC^2-(\overrightarrow{AB}\cdot\overrightarrow{AC})^2}$, $BC=\sqrt{53}$, đường cao $AH=\dfrac{2S}{BC}=\dfrac{\sqrt{901}}{53}$.}]{Cho tam giác $ABC$ có $A(1;-4;1)$, $B(-3;3;2)$, $C(1;-3;1)$. Độ dài đường cao kẻ từ $A$ của tam giác bằng}
  {$\dfrac{\sqrt{901}}{53}$.}
  {$\dfrac{\sqrt{3710}}{53}$.}{$\dfrac{\sqrt{1802}}{53}$.}{$\dfrac{\sqrt{1802}}{106}$.}

\cauTN[id=C2-B8-D10-TH-030,muc=TH,loigiai={Diện tích $S=\dfrac{1}{2}\sqrt{AB^2\cdot AC^2-(\overrightarrow{AB}\cdot\overrightarrow{AC})^2}$, $BC=\sqrt{51}$, đường cao $AH=\dfrac{2S}{BC}=\dfrac{\sqrt{100470}}{51}$.}]{Cho tam giác $ABC$ có $A(-3;-4;4)$, $B(-5;-3;-2)$, $C(-4;4;-3)$. Độ dài đường cao kẻ từ $A$ của tam giác bằng}
  {$\dfrac{\sqrt{100470}}{51}$.}
  {$\dfrac{\sqrt{103071}}{51}$.}{$\dfrac{2\sqrt{50235}}{51}$.}{$\dfrac{\sqrt{50235}}{51}$.}

\cauTN[id=C2-B8-D10-TH-031,muc=TH,loigiai={Diện tích $S=\dfrac{1}{2}\sqrt{AB^2\cdot AC^2-(\overrightarrow{AB}\cdot\overrightarrow{AC})^2}$, $BC=\sqrt{116}$, đường cao $AH=\dfrac{2S}{BC}=\dfrac{14\sqrt{145}}{29}$.}]{Cho tam giác $ABC$ có $A(-5;-2;1)$, $B(5;-4;5)$, $C(-1;0;-3)$. Độ dài đường cao kẻ từ $A$ của tam giác bằng}
  {$\dfrac{14\sqrt{145}}{29}$.}
  {$\dfrac{\sqrt{29261}}{29}$.}{$\dfrac{14\sqrt{290}}{29}$.}{$\dfrac{7\sqrt{290}}{29}$.}

\cauTN[id=C2-B8-D10-TH-032,muc=TH,loigiai={Diện tích $S=\dfrac{1}{2}\sqrt{AB^2\cdot AC^2-(\overrightarrow{AB}\cdot\overrightarrow{AC})^2}$, $BC=\sqrt{75}$, đường cao $AH=\dfrac{2S}{BC}=\dfrac{\sqrt{222}}{3}$.}]{Cho tam giác $ABC$ có $A(5;5;4)$, $B(4;5;-1)$, $C(-1;-2;-2)$. Độ dài đường cao kẻ từ $A$ của tam giác bằng}
  {$\dfrac{\sqrt{222}}{3}$.}
  {$\dfrac{\sqrt{231}}{3}$.}{$\dfrac{2\sqrt{111}}{3}$.}{$\dfrac{\sqrt{111}}{3}$.}

\cauTN[id=C2-B8-D10-TH-033,muc=TH,loigiai={Diện tích $S=\dfrac{1}{2}\sqrt{AB^2\cdot AC^2-(\overrightarrow{AB}\cdot\overrightarrow{AC})^2}$, $BC=\sqrt{35}$, đường cao $AH=\dfrac{2S}{BC}=\dfrac{\sqrt{2190}}{5}$.}]{Cho tam giác $ABC$ có $A(2;4;0)$, $B(-3;-4;0)$, $C(0;-5;5)$. Độ dài đường cao kẻ từ $A$ của tam giác bằng}
  {$\dfrac{\sqrt{2190}}{5}$.}
  {$\dfrac{\sqrt{2215}}{5}$.}{$\dfrac{2\sqrt{1095}}{5}$.}{$\dfrac{\sqrt{1095}}{5}$.}

\cauTN[id=C2-B8-D10-TH-034,muc=TH,loigiai={Diện tích $S=\dfrac{1}{2}\sqrt{AB^2\cdot AC^2-(\overrightarrow{AB}\cdot\overrightarrow{AC})^2}$, $BC=\sqrt{128}$, đường cao $AH=\dfrac{2S}{BC}=\dfrac{3\sqrt{2}}{2}$.}]{Cho tam giác $ABC$ có $A(-5;-2;-4)$, $B(3;3;-4)$, $C(-5;-5;-4)$. Độ dài đường cao kẻ từ $A$ của tam giác bằng}
  {$\dfrac{3\sqrt{2}}{2}$.}
  {$\dfrac{\sqrt{22}}{2}$.}{$3$.}{$\dfrac{3}{2}$.}

\cauTN[id=C2-B8-D10-TH-035,muc=TH,loigiai={Diện tích $S=\dfrac{1}{2}\sqrt{AB^2\cdot AC^2-(\overrightarrow{AB}\cdot\overrightarrow{AC})^2}$, $BC=\sqrt{153}$, đường cao $AH=\dfrac{2S}{BC}=\dfrac{3\sqrt{221}}{17}$.}]{Cho tam giác $ABC$ có $A(-5;1;5)$, $B(-5;-2;5)$, $C(4;4;-1)$. Độ dài đường cao kẻ từ $A$ của tam giác bằng}
  {$\dfrac{3\sqrt{221}}{17}$.}
  {$\dfrac{\sqrt{2278}}{17}$.}{$\dfrac{3\sqrt{442}}{17}$.}{$\dfrac{3\sqrt{442}}{34}$.}

\cauTN[id=C2-B8-D10-TH-036,muc=TH,loigiai={Diện tích $S=\dfrac{1}{2}\sqrt{AB^2\cdot AC^2-(\overrightarrow{AB}\cdot\overrightarrow{AC})^2}$, $BC=\sqrt{5}$, đường cao $AH=\dfrac{2S}{BC}=2\sqrt{5}$.}]{Cho tam giác $ABC$ có $A(-4;-5;3)$, $B(-4;3;2)$, $C(-4;5;3)$. Độ dài đường cao kẻ từ $A$ của tam giác bằng}
  {$2\sqrt{5}$.}
  {$\sqrt{21}$.}{$2\sqrt{10}$.}{$\sqrt{10}$.}

\cauTN[id=C2-B8-D10-TH-037,muc=TH,loigiai={Theo hệ trục đã chọn, $A(0;0;0)$, $S(0;0;5)$ nên $AS=5$.}]{Cho hình chóp $S.ABCD$ có đáy $ABCD$ là hình chữ nhật, $AB=5$, $AD=2$, $SA\perp(ABCD)$ và $SA=5$. Chọn hệ trục tọa độ $Oxyz$ có gốc $A$, tia $AB$ là $Ox$, tia $AD$ là $Oy$, tia $AS$ là $Oz$. Độ dài đoạn thẳng $AS$ bằng}
  {$5$.}
  {$\sqrt{26}$.}{$2\sqrt{6}$.}{$\sqrt{29}$.}

\cauTN[id=C2-B8-D10-TH-038,muc=TH,loigiai={Theo hệ trục đã chọn, $A(0;0;0)$, $B(4;0;0)$ nên $AB=4$.}]{Cho hình chóp $S.ABCD$ có đáy $ABCD$ là hình chữ nhật, $AB=4$, $AD=2$, $SA\perp(ABCD)$ và $SA=3$. Chọn hệ trục tọa độ $Oxyz$ có gốc $A$, tia $AB$ là $Ox$, tia $AD$ là $Oy$, tia $AS$ là $Oz$. Độ dài đoạn thẳng $AB$ bằng}
  {$4$.}
  {$\sqrt{17}$.}{$\sqrt{15}$.}{$2\sqrt{5}$.}

\cauTN[id=C2-B8-D10-TH-039,muc=TH,loigiai={Theo hệ trục đã chọn, $S(0;0;1)$, $C(3;1;0)$ nên $SC=\sqrt{11}$.}]{Cho hình chóp $S.ABCD$ có đáy $ABCD$ là hình chữ nhật, $AB=3$, $AD=1$, $SA\perp(ABCD)$ và $SA=1$. Chọn hệ trục tọa độ $Oxyz$ có gốc $A$, tia $AB$ là $Ox$, tia $AD$ là $Oy$, tia $AS$ là $Oz$. Độ dài đoạn thẳng $SC$ bằng}
  {$\sqrt{11}$.}
  {$2\sqrt{3}$.}{$\sqrt{10}$.}{$\sqrt{15}$.}

\cauTN[id=C2-B8-D10-TH-040,muc=TH,loigiai={Theo hệ trục đã chọn, $C(4;3;0)$, $A(0;0;0)$ nên $CA=5$.}]{Cho hình chóp $S.ABCD$ có đáy $ABCD$ là hình chữ nhật, $AB=4$, $AD=3$, $SA\perp(ABCD)$ và $SA=2$. Chọn hệ trục tọa độ $Oxyz$ có gốc $A$, tia $AB$ là $Ox$, tia $AD$ là $Oy$, tia $AS$ là $Oz$. Độ dài đoạn thẳng $CA$ bằng}
  {$5$.}
  {$\sqrt{26}$.}{$2\sqrt{6}$.}{$\sqrt{29}$.}

\cauTN[id=C2-B8-D10-TH-041,muc=TH,loigiai={Theo hệ trục đã chọn, $B(4;0;0)$, $D(0;4;0)$ nên $BD=4\sqrt{2}$.}]{Cho hình chóp $S.ABCD$ có đáy $ABCD$ là hình chữ nhật, $AB=4$, $AD=4$, $SA\perp(ABCD)$ và $SA=1$. Chọn hệ trục tọa độ $Oxyz$ có gốc $A$, tia $AB$ là $Ox$, tia $AD$ là $Oy$, tia $AS$ là $Oz$. Độ dài đoạn thẳng $BD$ bằng}
  {$4\sqrt{2}$.}
  {$\sqrt{33}$.}{$\sqrt{31}$.}{$6$.}

\cauTN[id=C2-B8-D10-TH-042,muc=TH,loigiai={Theo hệ trục đã chọn, $D(0;4;0)$, $A(0;0;0)$ nên $DA=4$.}]{Cho hình chóp $S.ABCD$ có đáy $ABCD$ là hình chữ nhật, $AB=4$, $AD=4$, $SA\perp(ABCD)$ và $SA=5$. Chọn hệ trục tọa độ $Oxyz$ có gốc $A$, tia $AB$ là $Ox$, tia $AD$ là $Oy$, tia $AS$ là $Oz$. Độ dài đoạn thẳng $DA$ bằng}
  {$4$.}
  {$\sqrt{17}$.}{$\sqrt{15}$.}{$2\sqrt{5}$.}

\cauTN[id=C2-B8-D10-TH-043,muc=TH,loigiai={Theo hệ trục đã chọn, $S(0;0;4)$, $D(0;1;0)$ nên $SD=\sqrt{17}$.}]{Cho hình chóp $S.ABCD$ có đáy $ABCD$ là hình chữ nhật, $AB=3$, $AD=1$, $SA\perp(ABCD)$ và $SA=4$. Chọn hệ trục tọa độ $Oxyz$ có gốc $A$, tia $AB$ là $Ox$, tia $AD$ là $Oy$, tia $AS$ là $Oz$. Độ dài đoạn thẳng $SD$ bằng}
  {$\sqrt{17}$.}
  {$3\sqrt{2}$.}{$4$.}{$\sqrt{21}$.}

\cauTN[id=C2-B8-D10-TH-044,muc=TH,loigiai={Theo hệ trục đã chọn, $A(0;0;0)$, $D(0;4;0)$ nên $AD=4$.}]{Cho hình chóp $S.ABCD$ có đáy $ABCD$ là hình chữ nhật, $AB=2$, $AD=4$, $SA\perp(ABCD)$ và $SA=5$. Chọn hệ trục tọa độ $Oxyz$ có gốc $A$, tia $AB$ là $Ox$, tia $AD$ là $Oy$, tia $AS$ là $Oz$. Độ dài đoạn thẳng $AD$ bằng}
  {$4$.}
  {$\sqrt{17}$.}{$\sqrt{15}$.}{$2\sqrt{5}$.}

\cauTN[id=C2-B8-D10-TH-045,muc=TH,loigiai={Theo hệ trục đã chọn, $C(2;4;0)$, $S(0;0;4)$ nên $CS=6$.}]{Cho hình chóp $S.ABCD$ có đáy $ABCD$ là hình chữ nhật, $AB=2$, $AD=4$, $SA\perp(ABCD)$ và $SA=4$. Chọn hệ trục tọa độ $Oxyz$ có gốc $A$, tia $AB$ là $Ox$, tia $AD$ là $Oy$, tia $AS$ là $Oz$. Độ dài đoạn thẳng $CS$ bằng}
  {$6$.}
  {$\sqrt{37}$.}{$\sqrt{35}$.}{$2\sqrt{10}$.}

\cauTN[id=C2-B8-D10-TH-046,muc=TH,loigiai={Theo hệ trục đã chọn, $A(0;0;0)$, $S(0;0;3)$ nên $AS=3$.}]{Cho hình chóp $S.ABCD$ có đáy $ABCD$ là hình chữ nhật, $AB=2$, $AD=1$, $SA\perp(ABCD)$ và $SA=3$. Chọn hệ trục tọa độ $Oxyz$ có gốc $A$, tia $AB$ là $Ox$, tia $AD$ là $Oy$, tia $AS$ là $Oz$. Độ dài đoạn thẳng $AS$ bằng}
  {$3$.}
  {$\sqrt{10}$.}{$2\sqrt{2}$.}{$\sqrt{13}$.}

\cauTN[id=C2-B8-D10-TH-047,muc=TH,loigiai={Theo hệ trục đã chọn, $C(2;1;0)$, $D(0;1;0)$ nên $CD=2$.}]{Cho hình chóp $S.ABCD$ có đáy $ABCD$ là hình chữ nhật, $AB=2$, $AD=1$, $SA\perp(ABCD)$ và $SA=3$. Chọn hệ trục tọa độ $Oxyz$ có gốc $A$, tia $AB$ là $Ox$, tia $AD$ là $Oy$, tia $AS$ là $Oz$. Độ dài đoạn thẳng $CD$ bằng}
  {$2$.}
  {$\sqrt{5}$.}{$\sqrt{3}$.}{$2\sqrt{2}$.}

\cauTN[id=C2-B8-D10-TH-048,muc=TH,loigiai={Theo hệ trục đã chọn, $C(5;5;0)$, $D(0;5;0)$ nên $CD=5$.}]{Cho hình chóp $S.ABCD$ có đáy $ABCD$ là hình chữ nhật, $AB=5$, $AD=5$, $SA\perp(ABCD)$ và $SA=4$. Chọn hệ trục tọa độ $Oxyz$ có gốc $A$, tia $AB$ là $Ox$, tia $AD$ là $Oy$, tia $AS$ là $Oz$. Độ dài đoạn thẳng $CD$ bằng}
  {$5$.}
  {$\sqrt{26}$.}{$2\sqrt{6}$.}{$\sqrt{29}$.}

\cauDS[id=C2-B8-D10-TH-049,muc=TH,loigiai={Kiểm tra từng mệnh đề bằng công thức tọa độ. a) Sai. b) Sai. c) Sai. d) Đúng.}]{Cho hình chóp $S.ABCD$ có đáy $ABCD$ là hình chữ nhật, $AB=4$, $AD=5$, $SA\perp(ABCD)$ và $SA=1$. Chọn hệ trục tọa độ $Oxyz$ có gốc $A$, tia $AB$ là $Ox$, tia $AD$ là $Oy$, tia $AS$ là $Oz$. Xét tính đúng sai của các mệnh đề sau:}{%
  \y{Tam giác $ABC$ có $A(3;-2;2)$, $B(-2;-4;1)$, $C(-5;-3;-4)$ thì đường cao kẻ từ $A$ bằng $\dfrac{\sqrt{26635}}{35}$.}{s}
  \y{Diện tích tam giác $ABC$ bằng $11\sqrt{3}$.}{s}
  \y{Độ dài đoạn thẳng $AC$ bằng $3\sqrt{5}$.}{s}
  \y{Tam giác $ABC$ có $A(3;-2;2)$, $B(-2;-4;1)$, $C(-5;-3;-4)$ thì đường cao kẻ từ $A$ bằng $\dfrac{11\sqrt{210}}{35}$.}{d}}

\cauDS[id=C2-B8-D10-TH-050,muc=TH,loigiai={Kiểm tra từng mệnh đề bằng công thức tọa độ. a) Đúng. b) Sai. c) Sai. d) Đúng.}]{Cho hình chóp $S.ABCD$ có đáy $ABCD$ là hình chữ nhật, $AB=1$, $AD=2$, $SA\perp(ABCD)$ và $SA=5$. Chọn hệ trục tọa độ $Oxyz$ có gốc $A$, tia $AB$ là $Ox$, tia $AD$ là $Oy$, tia $AS$ là $Oz$. Xét tính đúng sai của các mệnh đề sau:}{%
  \y{Tam giác $ABC$ có $A(0;-4;-3)$, $B(2;3;4)$, $C(0;-5;2)$ thì đường cao kẻ từ $A$ bằng $\dfrac{\sqrt{934}}{6}$.}{d}
  \y{Diện tích tam giác $ABC$ bằng $2\sqrt{467}$.}{s}
  \y{Tam giác $ABC$ có $A(0;-4;-3)$, $B(2;3;4)$, $C(0;-5;2)$ thì đường cao kẻ từ $A$ bằng $\dfrac{\sqrt{970}}{6}$.}{s}
  \y{Độ dài đoạn thẳng $CB$ bằng $2$.}{d}}

\cauDS[id=C2-B8-D10-TH-051,muc=TH,loigiai={Kiểm tra từng mệnh đề bằng công thức tọa độ. a) Sai. b) Sai. c) Đúng. d) Đúng.}]{Cho hình chóp $S.ABCD$ có đáy $ABCD$ là hình chữ nhật, $AB=3$, $AD=2$, $SA\perp(ABCD)$ và $SA=2$. Chọn hệ trục tọa độ $Oxyz$ có gốc $A$, tia $AB$ là $Ox$, tia $AD$ là $Oy$, tia $AS$ là $Oz$. Xét tính đúng sai của các mệnh đề sau:}{%
  \y{Tam giác $ABC$ có $A(0;2;3)$, $B(0;2;4)$, $C(0;2;8)$ có chu vi bằng $12$.}{s}
  \y{Tam giác $ABC$ có $A(-5;2;2)$, $B(1;-3;3)$, $C(1;1;-5)$ thì đường cao kẻ từ $A$ bằng $\dfrac{3\sqrt{290}}{10}$.}{s}
  \y{Độ dài đoạn thẳng $DA$ bằng $2$.}{d}
  \y{Tam giác $ABC$ có $A(-5;2;2)$, $B(1;-3;3)$, $C(1;1;-5)$ thì đường cao kẻ từ $A$ bằng $\dfrac{3\sqrt{145}}{5}$.}{d}}

\cauDS[id=C2-B8-D10-TH-052,muc=TH,loigiai={Kiểm tra từng mệnh đề bằng công thức tọa độ. a) Sai. b) Đúng. c) Đúng. d) Sai.}]{Cho hình chóp $S.ABCD$ có đáy $ABCD$ là hình chữ nhật, $AB=4$, $AD=2$, $SA\perp(ABCD)$ và $SA=4$. Chọn hệ trục tọa độ $Oxyz$ có gốc $A$, tia $AB$ là $Ox$, tia $AD$ là $Oy$, tia $AS$ là $Oz$. Xét tính đúng sai của các mệnh đề sau:}{%
  \y{Tam giác $ABC$ có $A(5;-1;-5)$, $B(-2;3;-2)$, $C(-1;5;3)$ thì đường cao kẻ từ $A$ bằng $\dfrac{\sqrt{14955}}{15}$.}{s}
  \y{Tam giác $ABC$ có $A(2;-2;-3)$, $B(0;0;-2)$, $C(4;4;0)$ có chu vi bằng $16$.}{d}
  \y{Diện tích tam giác $ABC$ bằng $\sqrt{491}$.}{d}
  \y{Độ dài đoạn thẳng $SB$ bằng $\sqrt{33}$.}{s}}

\cauDS[id=C2-B8-D10-TH-053,muc=TH,loigiai={Kiểm tra từng mệnh đề bằng công thức tọa độ. a) Đúng. b) Đúng. c) Sai. d) Đúng.}]{Cho hình chóp $S.ABCD$ có đáy $ABCD$ là hình chữ nhật, $AB=3$, $AD=1$, $SA\perp(ABCD)$ và $SA=5$. Chọn hệ trục tọa độ $Oxyz$ có gốc $A$, tia $AB$ là $Ox$, tia $AD$ là $Oy$, tia $AS$ là $Oz$. Xét tính đúng sai của các mệnh đề sau:}{%
  \y{Tam giác $ABC$ có $A(0;-1;0)$, $B(3;-1;-4)$, $C(2;-2;4)$ thì đường cao kẻ từ $A$ bằng $\dfrac{5\sqrt{1122}}{66}$.}{d}
  \y{Diện tích tam giác $ABC$ bằng $\dfrac{5\sqrt{17}}{2}$.}{d}
  \y{Tam giác $ABC$ có $A(0;-1;0)$, $B(3;-1;-4)$, $C(2;-2;4)$ thì đường cao kẻ từ $A$ bằng $\dfrac{5\sqrt{561}}{33}$.}{s}
  \y{Độ dài đoạn thẳng $DA$ bằng $1$.}{d}}

\cauDS[id=C2-B8-D10-TH-054,muc=TH,loigiai={Kiểm tra từng mệnh đề bằng công thức tọa độ. a) Sai. b) Đúng. c) Sai. d) Đúng.}]{Cho hình chóp $S.ABCD$ có đáy $ABCD$ là hình chữ nhật, $AB=1$, $AD=2$, $SA\perp(ABCD)$ và $SA=4$. Chọn hệ trục tọa độ $Oxyz$ có gốc $A$, tia $AB$ là $Ox$, tia $AD$ là $Oy$, tia $AS$ là $Oz$. Xét tính đúng sai của các mệnh đề sau:}{%
  \y{Diện tích tam giác $ABC$ bằng $\dfrac{\sqrt{922}}{2}$.}{s}
  \y{Tam giác $ABC$ có $A(-1;1;2)$, $B(2;-2;2)$, $C(2;-4;-5)$ thì đường cao kẻ từ $A$ bằng $\dfrac{3\sqrt{5406}}{53}$.}{d}
  \y{Diện tích tam giác $ABC$ bằng $3\sqrt{102}$.}{s}
  \y{Độ dài đoạn thẳng $DA$ bằng $2$.}{d}}

\cauDS[id=C2-B8-D10-TH-055,muc=TH,loigiai={Kiểm tra từng mệnh đề bằng công thức tọa độ. a) Sai. b) Sai. c) Sai. d) Đúng.}]{Cho hình chóp $S.ABCD$ có đáy $ABCD$ là hình chữ nhật, $AB=1$, $AD=4$, $SA\perp(ABCD)$ và $SA=2$. Chọn hệ trục tọa độ $Oxyz$ có gốc $A$, tia $AB$ là $Ox$, tia $AD$ là $Oy$, tia $AS$ là $Oz$. Xét tính đúng sai của các mệnh đề sau:}{%
  \y{Tam giác $ABC$ có $A(-4;4;2)$, $B(-1;-4;-4)$, $C(3;-4;3)$ thì đường cao kẻ từ $A$ bằng $\dfrac{\sqrt{32162}}{26}$.}{s}
  \y{Độ dài đoạn thẳng $CD$ bằng $\sqrt{3}$.}{s}
  \y{Diện tích tam giác $ABC$ bằng $\sqrt{6185}$.}{s}
  \y{Tam giác $ABC$ có $A(-2;-3;1)$, $B(-1;-3;1)$, $C(3;-3;1)$ có chu vi bằng $10$.}{d}}

\cauDS[id=C2-B8-D10-TH-056,muc=TH,loigiai={Kiểm tra từng mệnh đề bằng công thức tọa độ. a) Đúng. b) Sai. c) Sai. d) Đúng.}]{Trong không gian $Oxyz$, cho $A(-1;2;5)$, $B(1;-2;0)$, $C(-5;-2;-4)$. Xét tính đúng sai của các mệnh đề sau:}{%
  \y{Tam giác $ABC$ có $A(-1;2;5)$, $B(1;-2;0)$, $C(-5;-2;-4)$ thì đường cao kẻ từ $A$ bằng $\dfrac{\sqrt{7397}}{13}$.}{d}
  \y{Tam giác $ABC$ có $A(-2;3;3)$, $B(-2;8;3)$, $C(-2;4;3)$ có chu vi bằng $12$.}{s}
  \y{Tam giác $ABC$ có $A(1;3;2)$, $B(3;4;4)$, $C(3;2;4)$ có chu vi bằng $7$.}{s}
  \y{Diện tích tam giác $ABC$ bằng $\sqrt{569}$.}{d}}

\cauDS[id=C2-B8-D10-TH-057,muc=TH,loigiai={Kiểm tra từng mệnh đề bằng công thức tọa độ. a) Sai. b) Đúng. c) Đúng. d) Đúng.}]{Cho hình chóp $S.ABCD$ có đáy $ABCD$ là hình chữ nhật, $AB=4$, $AD=5$, $SA\perp(ABCD)$ và $SA=2$. Chọn hệ trục tọa độ $Oxyz$ có gốc $A$, tia $AB$ là $Ox$, tia $AD$ là $Oy$, tia $AS$ là $Oz$. Xét tính đúng sai của các mệnh đề sau:}{%
  \y{Tam giác $ABC$ có $A(-3;4;5)$, $B(-4;2;-5)$, $C(0;-1;-4)$ thì đường cao kẻ từ $A$ bằng $\dfrac{\sqrt{17498}}{13}$.}{s}
  \y{Độ dài đoạn thẳng $DB$ bằng $\sqrt{41}$.}{d}
  \y{Tam giác $ABC$ có $A(-3;4;5)$, $B(-4;2;-5)$, $C(0;-1;-4)$ thì đường cao kẻ từ $A$ bằng $\dfrac{\sqrt{17329}}{13}$.}{d}
  \y{Diện tích tam giác $ABC$ bằng $\dfrac{\sqrt{2666}}{2}$.}{d}}

\cauDS[id=C2-B8-D10-TH-058,muc=TH,loigiai={Kiểm tra từng mệnh đề bằng công thức tọa độ. a) Sai. b) Đúng. c) Đúng. d) Sai.}]{Cho hình chóp $S.ABCD$ có đáy $ABCD$ là hình chữ nhật, $AB=5$, $AD=3$, $SA\perp(ABCD)$ và $SA=5$. Chọn hệ trục tọa độ $Oxyz$ có gốc $A$, tia $AB$ là $Ox$, tia $AD$ là $Oy$, tia $AS$ là $Oz$. Xét tính đúng sai của các mệnh đề sau:}{%
  \y{Diện tích tam giác $ABC$ bằng $3\sqrt{199}$.}{s}
  \y{Diện tích tam giác $ABC$ bằng $\sqrt{1790}$.}{d}
  \y{Tam giác $ABC$ có $A(-3;1;1)$, $B(-3;1;2)$, $C(-3;1;6)$ có chu vi bằng $10$.}{d}
  \y{Độ dài đoạn thẳng $AS$ bằng $2\sqrt{6}$.}{s}}

\cauDS[id=C2-B8-D10-TH-059,muc=TH,loigiai={Kiểm tra từng mệnh đề bằng công thức tọa độ. a) Đúng. b) Sai. c) Sai. d) Sai.}]{Cho hình chóp $S.ABCD$ có đáy $ABCD$ là hình chữ nhật, $AB=4$, $AD=1$, $SA\perp(ABCD)$ và $SA=4$. Chọn hệ trục tọa độ $Oxyz$ có gốc $A$, tia $AB$ là $Ox$, tia $AD$ là $Oy$, tia $AS$ là $Oz$. Xét tính đúng sai của các mệnh đề sau:}{%
  \y{Độ dài đoạn thẳng $AS$ bằng $4$.}{d}
  \y{Diện tích tam giác $ABC$ bằng $\dfrac{3\sqrt{186}}{2}$.}{s}
  \y{Diện tích tam giác $ABC$ bằng $3\sqrt{93}$.}{s}
  \y{Tam giác $ABC$ có $A(2;0;-1)$, $B(3;2;1)$, $C(3;-2;1)$ có chu vi bằng $12$.}{s}}

\cauDS[id=C2-B8-D10-TH-060,muc=TH,loigiai={Kiểm tra từng mệnh đề bằng công thức tọa độ. a) Sai. b) Đúng. c) Sai. d) Sai.}]{Trong không gian $Oxyz$, cho $A(1;-1;-3)$, $B(2;-5;3)$, $C(-3;-4;-3)$. Xét tính đúng sai của các mệnh đề sau:}{%
  \y{Tam giác $ABC$ có $A(1;-1;-3)$, $B(2;-5;3)$, $C(-3;-4;-3)$ thì đường cao kẻ từ $A$ bằng $\dfrac{21\sqrt{186}}{62}$.}{s}
  \y{Diện tích tam giác $ABC$ bằng $\dfrac{\sqrt{1261}}{2}$.}{d}
  \y{Tam giác $ABC$ có $A(-2;-1;0)$, $B(-4;1;1)$, $C(0;1;1)$ có chu vi bằng $9$.}{s}
  \y{Diện tích tam giác $ABC$ bằng $\sqrt{1261}$.}{s}}

\cauDS[id=C2-B8-D10-TH-061,muc=TH,loigiai={Kiểm tra từng mệnh đề bằng công thức tọa độ. a) Sai. b) Sai. c) Đúng. d) Sai.}]{Cho hình chóp $S.ABCD$ có đáy $ABCD$ là hình chữ nhật, $AB=3$, $AD=2$, $SA\perp(ABCD)$ và $SA=2$. Chọn hệ trục tọa độ $Oxyz$ có gốc $A$, tia $AB$ là $Ox$, tia $AD$ là $Oy$, tia $AS$ là $Oz$. Xét tính đúng sai của các mệnh đề sau:}{%
  \y{Tam giác $ABC$ có $A(-3;-1;2)$, $B(-2;-3;4)$, $C(-2;1;4)$ có chu vi bằng $11$.}{s}
  \y{Diện tích tam giác $ABC$ bằng $\sqrt{2081}$.}{s}
  \y{Độ dài đoạn thẳng $CD$ bằng $3$.}{d}
  \y{Diện tích tam giác $ABC$ bằng $\dfrac{\sqrt{2085}}{2}$.}{s}}

\cauDS[id=C2-B8-D10-TH-062,muc=TH,loigiai={Kiểm tra từng mệnh đề bằng công thức tọa độ. a) Đúng. b) Sai. c) Sai. d) Sai.}]{Cho hình chóp $S.ABCD$ có đáy $ABCD$ là hình chữ nhật, $AB=3$, $AD=2$, $SA\perp(ABCD)$ và $SA=2$. Chọn hệ trục tọa độ $Oxyz$ có gốc $A$, tia $AB$ là $Ox$, tia $AD$ là $Oy$, tia $AS$ là $Oz$. Xét tính đúng sai của các mệnh đề sau:}{%
  \y{Diện tích tam giác $ABC$ bằng $\dfrac{\sqrt{1122}}{2}$.}{d}
  \y{Độ dài đoạn thẳng $SD$ bằng $3$.}{s}
  \y{Tam giác $ABC$ có $A(1;5;-1)$, $B(4;3;1)$, $C(-3;-2;-2)$ thì đường cao kẻ từ $A$ bằng $\dfrac{\sqrt{100015}}{83}$.}{s}
  \y{Tam giác $ABC$ có $A(-1;-2;2)$, $B(-1;-2;3)$, $C(-1;-2;7)$ có chu vi bằng $9$.}{s}}

\cauDS[id=C2-B8-D10-TH-063,muc=TH,loigiai={Kiểm tra từng mệnh đề bằng công thức tọa độ. a) Đúng. b) Sai. c) Đúng. d) Đúng.}]{Cho hình chóp $S.ABCD$ có đáy $ABCD$ là hình chữ nhật, $AB=3$, $AD=1$, $SA\perp(ABCD)$ và $SA=1$. Chọn hệ trục tọa độ $Oxyz$ có gốc $A$, tia $AB$ là $Ox$, tia $AD$ là $Oy$, tia $AS$ là $Oz$. Xét tính đúng sai của các mệnh đề sau:}{%
  \y{Tam giác $ABC$ có $A(0;2;-2)$, $B(-2;4;-1)$, $C(2;4;-1)$ có chu vi bằng $10$.}{d}
  \y{Diện tích tam giác $ABC$ bằng $2\sqrt{542}$.}{s}
  \y{Tam giác $ABC$ có $A(2;-1;0)$, $B(-4;-5;0)$, $C(-2;4;-1)$ thì đường cao kẻ từ $A$ bằng $\dfrac{2\sqrt{11653}}{43}$.}{d}
  \y{Độ dài đoạn thẳng $CB$ bằng $1$.}{d}}

\cauDS[id=C2-B8-D10-TH-064,muc=TH,loigiai={Kiểm tra từng mệnh đề bằng công thức tọa độ. a) Đúng. b) Đúng. c) Sai. d) Đúng.}]{Cho hình chóp $S.ABCD$ có đáy $ABCD$ là hình chữ nhật, $AB=5$, $AD=1$, $SA\perp(ABCD)$ và $SA=4$. Chọn hệ trục tọa độ $Oxyz$ có gốc $A$, tia $AB$ là $Ox$, tia $AD$ là $Oy$, tia $AS$ là $Oz$. Xét tính đúng sai của các mệnh đề sau:}{%
  \y{Tam giác $ABC$ có $A(0;5;0)$, $B(4;1;-3)$, $C(5;3;2)$ thì đường cao kẻ từ $A$ bằng $\dfrac{\sqrt{26070}}{30}$.}{d}
  \y{Diện tích tam giác $ABC$ bằng $\dfrac{\sqrt{869}}{2}$.}{d}
  \y{Diện tích tam giác $ABC$ bằng $\sqrt{869}$.}{s}
  \y{Độ dài đoạn thẳng $BC$ bằng $1$.}{d}}

\cauTLN[id=C2-B8-D10-TH-065,muc=TH,loigiai={$AB=5$, $AC=9$, $BC=6$ nên chu vi bằng $5+9+6=20$.}]{Tam giác $ABC$ có $A(-1;-1;0)$, $B(-1;-1;5)$, $C(3;3;7)$. Tính chu vi tam giác $ABC$.}{20}

\cauTLN[id=C2-B8-D10-TH-066,muc=TH,loigiai={$AB=5$, $AC=1$, $BC=4$ nên chu vi bằng $5+1+4=10$.}]{Tam giác $ABC$ có $A(0;2;3)$, $B(5;2;3)$, $C(1;2;3)$. Tính chu vi tam giác $ABC$.}{10}

\cauTLN[id=C2-B8-D10-TH-067,muc=TH,loigiai={$AB=9$, $AC=5$, $BC=6$ nên chu vi bằng $9+5+6=20$.}]{Tam giác $ABC$ có $A(3;-1;-2)$, $B(7;3;5)$, $C(3;-1;3)$. Tính chu vi tam giác $ABC$.}{20}

\cauTLN[id=C2-B8-D10-TH-068,muc=TH,loigiai={$AB=5$, $AC=9$, $BC=6$ nên chu vi bằng $5+9+6=20$.}]{Tam giác $ABC$ có $A(1;-1;-3)$, $B(1;-1;2)$, $C(5;3;4)$. Tính chu vi tam giác $ABC$.}{20}

\cauTLN[id=C2-B8-D10-TH-069,muc=TH,loigiai={$AB=1$, $AC=5$, $BC=4$ nên chu vi bằng $1+5+4=10$.}]{Tam giác $ABC$ có $A(-1;-3;-2)$, $B(-1;-3;-1)$, $C(-1;-3;3)$. Tính chu vi tam giác $ABC$.}{10}

\cauTLN[id=C2-B8-D10-TH-070,muc=TH,loigiai={$AB=1$, $AC=5$, $BC=4$ nên chu vi bằng $1+5+4=10$.}]{Tam giác $ABC$ có $A(1;-2;0)$, $B(1;-2;1)$, $C(1;-2;5)$. Tính chu vi tam giác $ABC$.}{10}

\cauTLN[id=C2-B8-D10-TH-071,muc=TH,loigiai={Theo hệ trục đã chọn, $A(0;0;0)$, $S(0;0;1)$ nên $AS=1$.}]{Cho hình chóp $S.ABCD$ có đáy $ABCD$ là hình chữ nhật, $AB=5$, $AD=2$, $SA\perp(ABCD)$ và $SA=1$. Chọn hệ trục tọa độ $Oxyz$ có gốc $A$, tia $AB$ là $Ox$, tia $AD$ là $Oy$, tia $AS$ là $Oz$. Tính $AS^2$.}{1}

\cauTLN[id=C2-B8-D10-TH-072,muc=TH,loigiai={Theo hệ trục đã chọn, $B(1;0;0)$, $S(0;0;4)$ nên $BS=\sqrt{17}$.}]{Cho hình chóp $S.ABCD$ có đáy $ABCD$ là hình chữ nhật, $AB=1$, $AD=5$, $SA\perp(ABCD)$ và $SA=4$. Chọn hệ trục tọa độ $Oxyz$ có gốc $A$, tia $AB$ là $Ox$, tia $AD$ là $Oy$, tia $AS$ là $Oz$. Tính $BS^2$.}{17}

\cauTLN[id=C2-B8-D10-TH-073,muc=TH,loigiai={Theo hệ trục đã chọn, $D(0;3;0)$, $B(3;0;0)$ nên $DB=3\sqrt{2}$.}]{Cho hình chóp $S.ABCD$ có đáy $ABCD$ là hình chữ nhật, $AB=3$, $AD=3$, $SA\perp(ABCD)$ và $SA=1$. Chọn hệ trục tọa độ $Oxyz$ có gốc $A$, tia $AB$ là $Ox$, tia $AD$ là $Oy$, tia $AS$ là $Oz$. Tính $DB^2$.}{18}

\cauTLN[id=C2-B8-D10-TH-074,muc=TH,loigiai={Theo hệ trục đã chọn, $S(0;0;1)$, $D(0;3;0)$ nên $SD=\sqrt{10}$.}]{Cho hình chóp $S.ABCD$ có đáy $ABCD$ là hình chữ nhật, $AB=2$, $AD=3$, $SA\perp(ABCD)$ và $SA=1$. Chọn hệ trục tọa độ $Oxyz$ có gốc $A$, tia $AB$ là $Ox$, tia $AD$ là $Oy$, tia $AS$ là $Oz$. Tính $SD^2$.}{10}

\cauTLN[id=C2-B8-D10-TH-075,muc=TH,loigiai={Theo hệ trục đã chọn, $S(0;0;5)$, $A(0;0;0)$ nên $SA=5$.}]{Cho hình chóp $S.ABCD$ có đáy $ABCD$ là hình chữ nhật, $AB=4$, $AD=5$, $SA\perp(ABCD)$ và $SA=5$. Chọn hệ trục tọa độ $Oxyz$ có gốc $A$, tia $AB$ là $Ox$, tia $AD$ là $Oy$, tia $AS$ là $Oz$. Tính $SA^2$.}{25}

\cauTLN[id=C2-B8-D10-TH-076,muc=TH,loigiai={Ta có $\overrightarrow{AB}=(-4;-1;5)$, $\overrightarrow{AC}=(-7;3;-3)$. Diện tích $S=\dfrac{1}{2}\sqrt{|\overrightarrow{AB}|^2|\overrightarrow{AC}|^2-(\overrightarrow{AB}\cdot\overrightarrow{AC})^2}=\dfrac{1}{2}\sqrt{42\cdot67-10^2}=\dfrac{\sqrt{2714}}{2}$.}]{Trong không gian $Oxyz$, cho $A(4;0;0)$, $B(0;-1;5)$, $C(-3;3;-3)$. Tính $4S^2$, trong đó $S$ là diện tích tam giác $ABC$.}{2714}

\cauTLN[id=C2-B8-D10-TH-077,muc=TH,loigiai={Ta có $\overrightarrow{AB}=(6;2;-7)$, $\overrightarrow{AC}=(4;-1;-3)$. Diện tích $S=\dfrac{1}{2}\sqrt{|\overrightarrow{AB}|^2|\overrightarrow{AC}|^2-(\overrightarrow{AB}\cdot\overrightarrow{AC})^2}=\dfrac{1}{2}\sqrt{89\cdot26-43^2}=\dfrac{\sqrt{465}}{2}$.}]{Trong không gian $Oxyz$, cho $A(-3;2;2)$, $B(3;4;-5)$, $C(1;1;-1)$. Tính $4S^2$, trong đó $S$ là diện tích tam giác $ABC$.}{465}

\cauTLN[id=C2-B8-D10-TH-078,muc=TH,loigiai={Ta có $\overrightarrow{AB}=(-3;-6;1)$, $\overrightarrow{AC}=(2;-1;1)$. Diện tích $S=\dfrac{1}{2}\sqrt{|\overrightarrow{AB}|^2|\overrightarrow{AC}|^2-(\overrightarrow{AB}\cdot\overrightarrow{AC})^2}=\dfrac{1}{2}\sqrt{46\cdot6-1^2}=\dfrac{5\sqrt{11}}{2}$.}]{Trong không gian $Oxyz$, cho $A(2;5;3)$, $B(-1;-1;4)$, $C(4;4;4)$. Tính $4S^2$, trong đó $S$ là diện tích tam giác $ABC$.}{275}

\cauTLN[id=C2-B8-D10-TH-079,muc=TH,loigiai={Ta có $\overrightarrow{AB}=(-1;6;-5)$, $\overrightarrow{AC}=(-4;9;-5)$. Diện tích $S=\dfrac{1}{2}\sqrt{|\overrightarrow{AB}|^2|\overrightarrow{AC}|^2-(\overrightarrow{AB}\cdot\overrightarrow{AC})^2}=\dfrac{1}{2}\sqrt{62\cdot122-83^2}=\dfrac{15\sqrt{3}}{2}$.}]{Trong không gian $Oxyz$, cho $A(3;-4;2)$, $B(2;2;-3)$, $C(-1;5;-3)$. Tính $4S^2$, trong đó $S$ là diện tích tam giác $ABC$.}{675}

\cauTLN[id=C2-B8-D10-TH-080,muc=TH,loigiai={Ta có $\overrightarrow{AB}=(-4;-6;-2)$, $\overrightarrow{AC}=(-5;1;-4)$. Diện tích $S=\dfrac{1}{2}\sqrt{|\overrightarrow{AB}|^2|\overrightarrow{AC}|^2-(\overrightarrow{AB}\cdot\overrightarrow{AC})^2}=\dfrac{1}{2}\sqrt{56\cdot42-22^2}=\sqrt{467}$.}]{Trong không gian $Oxyz$, cho $A(2;1;4)$, $B(-2;-5;2)$, $C(-3;2;0)$. Tính $4S^2$, trong đó $S$ là diện tích tam giác $ABC$.}{1868}

\cauTN[id=C2-B8-D10-VD-001,muc=VD,loigiai={Gọi $M(x;y;0)$. Từ $MA^2=MB^2$ và $MA^2=MC^2$ ta được hệ hai phương trình bậc nhất ẩn $x,y$, giải được $x=0$, $y=2$. Vậy $M(0;2;0)$.}]{Cho $A(2;1;2)$, $B(1;0;2)$, $C(-2;4;1)$. Điểm $M$ thuộc mặt phẳng $(Oxy)$ cách đều ba điểm $A,B,C$ có tọa độ là}
  {$(0;2;0)$.}
  {$(2;2;0)$.}{$(-1;2;0)$.}{$(2;0;0)$.}

\cauTN[id=C2-B8-D10-VD-002,muc=VD,loigiai={Gọi $M(x;y;0)$. Từ $MA^2=MB^2$ và $MA^2=MC^2$ ta được hệ hai phương trình bậc nhất ẩn $x,y$, giải được $x=3$, $y=-2$. Vậy $M(3;-2;0)$.}]{Cho $A(5;-3;2)$, $B(5;0;-1)$, $C(4;0;2)$. Điểm $M$ thuộc mặt phẳng $(Oxy)$ cách đều ba điểm $A,B,C$ có tọa độ là}
  {$(3;-2;0)$.}
  {$(3;0;-2)$.}{$(5;-2;0)$.}{$(-2;3;0)$.}

\cauTN[id=C2-B8-D10-VD-003,muc=VD,loigiai={Gọi $M(x;y;0)$. Từ $MA^2=MB^2$ và $MA^2=MC^2$ ta được hệ hai phương trình bậc nhất ẩn $x,y$, giải được $x=-2$, $y=-1$. Vậy $M(-2;-1;0)$.}]{Cho $A(0;1;1)$, $B(-4;1;1)$, $C(-1;-3;2)$. Điểm $M$ thuộc mặt phẳng $(Oxy)$ cách đều ba điểm $A,B,C$ có tọa độ là}
  {$(-2;-1;0)$.}
  {$(-1;-1;0)$.}{$(-2;-1;-1)$.}{$(-2;-3;0)$.}

\cauTN[id=C2-B8-D10-VD-004,muc=VD,loigiai={Gọi $M(x;y;0)$. Từ $MA^2=MB^2$ và $MA^2=MC^2$ ta được hệ hai phương trình bậc nhất ẩn $x,y$, giải được $x=-2$, $y=2$. Vậy $M(-2;2;0)$.}]{Cho $A(0;5;6)$, $B(4;5;2)$, $C(0;8;3)$. Điểm $M$ thuộc mặt phẳng $(Oxy)$ cách đều ba điểm $A,B,C$ có tọa độ là}
  {$(-2;2;0)$.}
  {$(2;2;0)$.}{$(-2;1;0)$.}{$(-2;0;2)$.}

\cauTN[id=C2-B8-D10-VD-005,muc=VD,loigiai={Gọi $M(x;y;0)$. Từ $MA^2=MB^2$ và $MA^2=MC^2$ ta được hệ hai phương trình bậc nhất ẩn $x,y$, giải được $x=3$, $y=2$. Vậy $M(3;2;0)$.}]{Cho $A(3;2;5)$, $B(7;5;0)$, $C(3;7;0)$. Điểm $M$ thuộc mặt phẳng $(Oxy)$ cách đều ba điểm $A,B,C$ có tọa độ là}
  {$(3;2;0)$.}
  {$(-3;-2;0)$.}{$(3;0;2)$.}{$(3;4;0)$.}

\cauTN[id=C2-B8-D10-VD-006,muc=VD,loigiai={Gọi $M(x;y;0)$. Từ $MA^2=MB^2$ và $MA^2=MC^2$ ta được hệ hai phương trình bậc nhất ẩn $x,y$, giải được $x=0$, $y=-3$. Vậy $M(0;-3;0)$.}]{Cho $A(0;1;3)$, $B(0;0;4)$, $C(5;-3;0)$. Điểm $M$ thuộc mặt phẳng $(Oxy)$ cách đều ba điểm $A,B,C$ có tọa độ là}
  {$(0;-3;0)$.}
  {$(0;-3;-2)$.}{$(0;0;-3)$.}{$(-2;-3;0)$.}

\cauTN[id=C2-B8-D10-VD-007,muc=VD,loigiai={Gọi $M(x;y;0)$. Từ $MA^2=MB^2$ và $MA^2=MC^2$ ta được hệ hai phương trình bậc nhất ẩn $x,y$, giải được $x=0$, $y=3$. Vậy $M(0;3;0)$.}]{Cho $A(0;6;4)$, $B(4;6;0)$, $C(3;3;4)$. Điểm $M$ thuộc mặt phẳng $(Oxy)$ cách đều ba điểm $A,B,C$ có tọa độ là}
  {$(0;3;0)$.}
  {$(0;-3;0)$.}{$(0;3;2)$.}{$(1;3;0)$.}

\cauTN[id=C2-B8-D10-VD-008,muc=VD,loigiai={Gọi $M(x;y;0)$. Từ $MA^2=MB^2$ và $MA^2=MC^2$ ta được hệ hai phương trình bậc nhất ẩn $x,y$, giải được $x=1$, $y=2$. Vậy $M(1;2;0)$.}]{Cho $A(2;2;0)$, $B(1;2;1)$, $C(1;3;0)$. Điểm $M$ thuộc mặt phẳng $(Oxy)$ cách đều ba điểm $A,B,C$ có tọa độ là}
  {$(1;2;0)$.}
  {$(1;4;0)$.}{$(1;2;-1)$.}{$(-1;2;0)$.}

\cauTN[id=C2-B8-D10-VD-009,muc=VD,loigiai={Gọi $M(x;y;0)$. Từ $MA^2=MB^2$ và $MA^2=MC^2$ ta được hệ hai phương trình bậc nhất ẩn $x,y$, giải được $x=3$, $y=-3$. Vậy $M(3;-3;0)$.}]{Cho $A(9;0;2)$, $B(6;-1;6)$, $C(5;0;6)$. Điểm $M$ thuộc mặt phẳng $(Oxy)$ cách đều ba điểm $A,B,C$ có tọa độ là}
  {$(3;-3;0)$.}
  {$(3;-3;-2)$.}{$(-3;3;0)$.}{$(2;-3;0)$.}

\cauTN[id=C2-B8-D10-VD-010,muc=VD,loigiai={Gọi $M(x;y;0)$. Từ $MA^2=MB^2$ và $MA^2=MC^2$ ta được hệ hai phương trình bậc nhất ẩn $x,y$, giải được $x=2$, $y=0$. Vậy $M(2;0;0)$.}]{Cho $A(2;3;4)$, $B(6;0;3)$, $C(2;4;3)$. Điểm $M$ thuộc mặt phẳng $(Oxy)$ cách đều ba điểm $A,B,C$ có tọa độ là}
  {$(2;0;0)$.}
  {$(0;0;0)$.}{$(3;0;0)$.}{$(4;0;0)$.}

\cauTN[id=C2-B8-D10-VD-011,muc=VD,loigiai={Gọi $M(x;y;0)$. Từ $MA^2=MB^2$ và $MA^2=MC^2$ ta được hệ hai phương trình bậc nhất ẩn $x,y$, giải được $x=2$, $y=-2$. Vậy $M(2;-2;0)$.}]{Cho $A(6;-2;3)$, $B(2;-2;5)$, $C(6;1;0)$. Điểm $M$ thuộc mặt phẳng $(Oxy)$ cách đều ba điểm $A,B,C$ có tọa độ là}
  {$(2;-2;0)$.}
  {$(-2;2;0)$.}{$(1;-2;0)$.}{$(2;-2;-1)$.}

\cauTN[id=C2-B8-D10-VD-012,muc=VD,loigiai={Gọi $M(x;y;0)$. Từ $MA^2=MB^2$ và $MA^2=MC^2$ ta được hệ hai phương trình bậc nhất ẩn $x,y$, giải được $x=3$, $y=-1$. Vậy $M(3;-1;0)$.}]{Cho $A(1;1;1)$, $B(5;-2;2)$, $C(5;0;2)$. Điểm $M$ thuộc mặt phẳng $(Oxy)$ cách đều ba điểm $A,B,C$ có tọa độ là}
  {$(3;-1;0)$.}
  {$(0;-1;3)$.}{$(3;-2;0)$.}{$(4;-1;0)$.}

\cauTN[id=C2-B8-D10-VD-013,muc=VD,loigiai={Có $\overrightarrow{SC}=(1;3;-2)$, $\overrightarrow{BD}=(-1;3;0)$ nên $\cos(\overrightarrow{SC},\overrightarrow{BD})=\dfrac{4\sqrt{35}}{35}$.}]{Cho hình chóp $S.ABCD$ có đáy $ABCD$ là hình chữ nhật, $AB=1$, $AD=3$, $SA\perp(ABCD)$ và $SA=2$. Chọn hệ trục tọa độ $Oxyz$ có gốc $A$, tia $AB$ là $Ox$, tia $AD$ là $Oy$, tia $AS$ là $Oz$. Giá trị $\cos(\overrightarrow{SC},\overrightarrow{BD})$ bằng}
  {$\dfrac{4\sqrt{35}}{35}$.}
  {$-\dfrac{4\sqrt{35}}{35}$.}{$\dfrac{9\sqrt{35}}{70}$.}{$\dfrac{\sqrt{35}}{10}$.}

\cauTN[id=C2-B8-D10-VD-014,muc=VD,loigiai={Có $\overrightarrow{SC}=(1;5;-3)$, $\overrightarrow{BD}=(-1;5;0)$ nên $\cos(\overrightarrow{SC},\overrightarrow{BD})=\dfrac{12\sqrt{910}}{455}$.}]{Cho hình chóp $S.ABCD$ có đáy $ABCD$ là hình chữ nhật, $AB=1$, $AD=5$, $SA\perp(ABCD)$ và $SA=3$. Chọn hệ trục tọa độ $Oxyz$ có gốc $A$, tia $AB$ là $Ox$, tia $AD$ là $Oy$, tia $AS$ là $Oz$. Giá trị $\cos(\overrightarrow{SC},\overrightarrow{BD})$ bằng}
  {$\dfrac{12\sqrt{910}}{455}$.}
  {$-\dfrac{12\sqrt{910}}{455}$.}{$\dfrac{5\sqrt{910}}{182}$.}{$\dfrac{23\sqrt{910}}{910}$.}

\cauTN[id=C2-B8-D10-VD-015,muc=VD,loigiai={Có $\overrightarrow{SC}=(1;2;-4)$, $\overrightarrow{BD}=(-1;2;0)$ nên $\cos(\overrightarrow{SC},\overrightarrow{BD})=\dfrac{\sqrt{105}}{35}$.}]{Cho hình chóp $S.ABCD$ có đáy $ABCD$ là hình chữ nhật, $AB=1$, $AD=2$, $SA\perp(ABCD)$ và $SA=4$. Chọn hệ trục tọa độ $Oxyz$ có gốc $A$, tia $AB$ là $Ox$, tia $AD$ là $Oy$, tia $AS$ là $Oz$. Giá trị $\cos(\overrightarrow{SC},\overrightarrow{BD})$ bằng}
  {$\dfrac{\sqrt{105}}{35}$.}
  {$-\dfrac{\sqrt{105}}{35}$.}{$\dfrac{4\sqrt{105}}{105}$.}{$\dfrac{2\sqrt{105}}{105}$.}

\cauTN[id=C2-B8-D10-VD-016,muc=VD,loigiai={Có $\overrightarrow{SC}=(3;1;-4)$, $\overrightarrow{BD}=(-3;1;0)$ nên $\cos(\overrightarrow{SC},\overrightarrow{BD})=-\dfrac{4\sqrt{65}}{65}$.}]{Cho hình chóp $S.ABCD$ có đáy $ABCD$ là hình chữ nhật, $AB=3$, $AD=1$, $SA\perp(ABCD)$ và $SA=4$. Chọn hệ trục tọa độ $Oxyz$ có gốc $A$, tia $AB$ là $Ox$, tia $AD$ là $Oy$, tia $AS$ là $Oz$. Giá trị $\cos(\overrightarrow{SC},\overrightarrow{BD})$ bằng}
  {$-\dfrac{4\sqrt{65}}{65}$.}
  {$\dfrac{4\sqrt{65}}{65}$.}{$-\dfrac{7\sqrt{65}}{130}$.}{$-\dfrac{9\sqrt{65}}{130}$.}

\cauTN[id=C2-B8-D10-VD-017,muc=VD,loigiai={Có $\overrightarrow{SC}=(2;1;-5)$, $\overrightarrow{BD}=(-2;1;0)$ nên $\cos(\overrightarrow{SC},\overrightarrow{BD})=-\dfrac{\sqrt{6}}{10}$.}]{Cho hình chóp $S.ABCD$ có đáy $ABCD$ là hình chữ nhật, $AB=2$, $AD=1$, $SA\perp(ABCD)$ và $SA=5$. Chọn hệ trục tọa độ $Oxyz$ có gốc $A$, tia $AB$ là $Ox$, tia $AD$ là $Oy$, tia $AS$ là $Oz$. Giá trị $\cos(\overrightarrow{SC},\overrightarrow{BD})$ bằng}
  {$-\dfrac{\sqrt{6}}{10}$.}
  {$\dfrac{\sqrt{6}}{10}$.}{$-\dfrac{\sqrt{6}}{15}$.}{$-\dfrac{2\sqrt{6}}{15}$.}

\cauTN[id=C2-B8-D10-VD-018,muc=VD,loigiai={Có $\overrightarrow{SC}=(3;1;-3)$, $\overrightarrow{BD}=(-3;1;0)$ nên $\cos(\overrightarrow{SC},\overrightarrow{BD})=-\dfrac{4\sqrt{190}}{95}$.}]{Cho hình chóp $S.ABCD$ có đáy $ABCD$ là hình chữ nhật, $AB=3$, $AD=1$, $SA\perp(ABCD)$ và $SA=3$. Chọn hệ trục tọa độ $Oxyz$ có gốc $A$, tia $AB$ là $Ox$, tia $AD$ là $Oy$, tia $AS$ là $Oz$. Giá trị $\cos(\overrightarrow{SC},\overrightarrow{BD})$ bằng}
  {$-\dfrac{4\sqrt{190}}{95}$.}
  {$\dfrac{4\sqrt{190}}{95}$.}{$-\dfrac{7\sqrt{190}}{190}$.}{$-\dfrac{9\sqrt{190}}{190}$.}

\cauTN[id=C2-B8-D10-VD-019,muc=VD,loigiai={Có $\overrightarrow{SC}=(5;4;-1)$, $\overrightarrow{BD}=(-5;4;0)$ nên $\cos(\overrightarrow{SC},\overrightarrow{BD})=-\dfrac{3\sqrt{1722}}{574}$.}]{Cho hình chóp $S.ABCD$ có đáy $ABCD$ là hình chữ nhật, $AB=5$, $AD=4$, $SA\perp(ABCD)$ và $SA=1$. Chọn hệ trục tọa độ $Oxyz$ có gốc $A$, tia $AB$ là $Ox$, tia $AD$ là $Oy$, tia $AS$ là $Oz$. Giá trị $\cos(\overrightarrow{SC},\overrightarrow{BD})$ bằng}
  {$-\dfrac{3\sqrt{1722}}{574}$.}
  {$\dfrac{3\sqrt{1722}}{574}$.}{$-\dfrac{4\sqrt{1722}}{861}$.}{$-\dfrac{5\sqrt{1722}}{861}$.}

\cauTN[id=C2-B8-D10-VD-020,muc=VD,loigiai={Có $\overrightarrow{SC}=(1;3;-1)$, $\overrightarrow{BD}=(-1;3;0)$ nên $\cos(\overrightarrow{SC},\overrightarrow{BD})=\dfrac{4\sqrt{110}}{55}$.}]{Cho hình chóp $S.ABCD$ có đáy $ABCD$ là hình chữ nhật, $AB=1$, $AD=3$, $SA\perp(ABCD)$ và $SA=1$. Chọn hệ trục tọa độ $Oxyz$ có gốc $A$, tia $AB$ là $Ox$, tia $AD$ là $Oy$, tia $AS$ là $Oz$. Giá trị $\cos(\overrightarrow{SC},\overrightarrow{BD})$ bằng}
  {$\dfrac{4\sqrt{110}}{55}$.}
  {$-\dfrac{4\sqrt{110}}{55}$.}{$\dfrac{9\sqrt{110}}{110}$.}{$\dfrac{7\sqrt{110}}{110}$.}

\cauTN[id=C2-B8-D10-VD-021,muc=VD,loigiai={Có $\overrightarrow{SC}=(4;5;-5)$, $\overrightarrow{BD}=(-4;5;0)$ nên $\cos(\overrightarrow{SC},\overrightarrow{BD})=\dfrac{3\sqrt{2706}}{902}$.}]{Cho hình chóp $S.ABCD$ có đáy $ABCD$ là hình chữ nhật, $AB=4$, $AD=5$, $SA\perp(ABCD)$ và $SA=5$. Chọn hệ trục tọa độ $Oxyz$ có gốc $A$, tia $AB$ là $Ox$, tia $AD$ là $Oy$, tia $AS$ là $Oz$. Giá trị $\cos(\overrightarrow{SC},\overrightarrow{BD})$ bằng}
  {$\dfrac{3\sqrt{2706}}{902}$.}
  {$-\dfrac{3\sqrt{2706}}{902}$.}{$\dfrac{5\sqrt{2706}}{1353}$.}{$\dfrac{4\sqrt{2706}}{1353}$.}

\cauTN[id=C2-B8-D10-VD-022,muc=VD,loigiai={Có $\overrightarrow{SC}=(4;1;-3)$, $\overrightarrow{BD}=(-4;1;0)$ nên $\cos(\overrightarrow{SC},\overrightarrow{BD})=-\dfrac{15\sqrt{442}}{442}$.}]{Cho hình chóp $S.ABCD$ có đáy $ABCD$ là hình chữ nhật, $AB=4$, $AD=1$, $SA\perp(ABCD)$ và $SA=3$. Chọn hệ trục tọa độ $Oxyz$ có gốc $A$, tia $AB$ là $Ox$, tia $AD$ là $Oy$, tia $AS$ là $Oz$. Giá trị $\cos(\overrightarrow{SC},\overrightarrow{BD})$ bằng}
  {$-\dfrac{15\sqrt{442}}{442}$.}
  {$\dfrac{15\sqrt{442}}{442}$.}{$-\dfrac{7\sqrt{442}}{221}$.}{$-\dfrac{8\sqrt{442}}{221}$.}

\cauTN[id=C2-B8-D10-VD-023,muc=VD,loigiai={Có $\overrightarrow{SC}=(3;1;-1)$, $\overrightarrow{BD}=(-3;1;0)$ nên $\cos(\overrightarrow{SC},\overrightarrow{BD})=-\dfrac{4\sqrt{110}}{55}$.}]{Cho hình chóp $S.ABCD$ có đáy $ABCD$ là hình chữ nhật, $AB=3$, $AD=1$, $SA\perp(ABCD)$ và $SA=1$. Chọn hệ trục tọa độ $Oxyz$ có gốc $A$, tia $AB$ là $Ox$, tia $AD$ là $Oy$, tia $AS$ là $Oz$. Giá trị $\cos(\overrightarrow{SC},\overrightarrow{BD})$ bằng}
  {$-\dfrac{4\sqrt{110}}{55}$.}
  {$\dfrac{4\sqrt{110}}{55}$.}{$-\dfrac{7\sqrt{110}}{110}$.}{$-\dfrac{9\sqrt{110}}{110}$.}

\cauTN[id=C2-B8-D10-VD-024,muc=VD,loigiai={Có $\overrightarrow{SC}=(5;1;-1)$, $\overrightarrow{BD}=(-5;1;0)$ nên $\cos(\overrightarrow{SC},\overrightarrow{BD})=-\dfrac{4\sqrt{78}}{39}$.}]{Cho hình chóp $S.ABCD$ có đáy $ABCD$ là hình chữ nhật, $AB=5$, $AD=1$, $SA\perp(ABCD)$ và $SA=1$. Chọn hệ trục tọa độ $Oxyz$ có gốc $A$, tia $AB$ là $Ox$, tia $AD$ là $Oy$, tia $AS$ là $Oz$. Giá trị $\cos(\overrightarrow{SC},\overrightarrow{BD})$ bằng}
  {$-\dfrac{4\sqrt{78}}{39}$.}
  {$\dfrac{4\sqrt{78}}{39}$.}{$-\dfrac{23\sqrt{78}}{234}$.}{$-\dfrac{25\sqrt{78}}{234}$.}

\cauTN[id=C2-B8-D10-VD-025,muc=VD,loigiai={Diện tích $S=\dfrac{1}{2}\sqrt{AB^2\cdot AC^2-(\overrightarrow{AB}\cdot\overrightarrow{AC})^2}$, $BC=\sqrt{4}$, đường cao $AH=\dfrac{2S}{BC}=3\sqrt{2}$.}]{Cho tam giác $ABC$ có $A(-1;0;-3)$, $B(-4;-3;-6)$, $C(-6;-3;-6)$. Độ dài đường cao kẻ từ $A$ của tam giác bằng}
  {$3\sqrt{2}$.}
  {$\sqrt{19}$.}{$6$.}{$3$.}

\cauTN[id=C2-B8-D10-VD-026,muc=VD,loigiai={Diện tích $S=\dfrac{1}{2}\sqrt{AB^2\cdot AC^2-(\overrightarrow{AB}\cdot\overrightarrow{AC})^2}$, $BC=\sqrt{225}$, đường cao $AH=\dfrac{2S}{BC}=\dfrac{41}{5}$.}]{Cho tam giác $ABC$ có $A(4;5;6)$, $B(4;-6;5)$, $C(2;5;-5)$. Độ dài đường cao kẻ từ $A$ của tam giác bằng}
  {$\dfrac{41}{5}$.}
  {$\dfrac{\sqrt{1706}}{5}$.}{$\dfrac{41\sqrt{2}}{5}$.}{$\dfrac{41\sqrt{2}}{10}$.}

\cauTN[id=C2-B8-D10-VD-027,muc=VD,loigiai={Diện tích $S=\dfrac{1}{2}\sqrt{AB^2\cdot AC^2-(\overrightarrow{AB}\cdot\overrightarrow{AC})^2}$, $BC=\sqrt{56}$, đường cao $AH=\dfrac{2S}{BC}=\dfrac{3\sqrt{70}}{14}$.}]{Cho tam giác $ABC$ có $A(1;0;5)$, $B(5;-4;-2)$, $C(3;0;4)$. Độ dài đường cao kẻ từ $A$ của tam giác bằng}
  {$\dfrac{3\sqrt{70}}{14}$.}
  {$\dfrac{\sqrt{826}}{14}$.}{$\dfrac{3\sqrt{35}}{7}$.}{$\dfrac{3\sqrt{35}}{14}$.}

\cauTN[id=C2-B8-D10-VD-028,muc=VD,loigiai={Diện tích $S=\dfrac{1}{2}\sqrt{AB^2\cdot AC^2-(\overrightarrow{AB}\cdot\overrightarrow{AC})^2}$, $BC=\sqrt{110}$, đường cao $AH=\dfrac{2S}{BC}=\dfrac{\sqrt{173085}}{55}$.}]{Cho tam giác $ABC$ có $A(0;3;3)$, $B(2;-1;-4)$, $C(-5;-6;2)$. Độ dài đường cao kẻ từ $A$ của tam giác bằng}
  {$\dfrac{\sqrt{173085}}{55}$.}
  {$\dfrac{\sqrt{176110}}{55}$.}{$\dfrac{\sqrt{346170}}{55}$.}{$\dfrac{\sqrt{346170}}{110}$.}

\cauTN[id=C2-B8-D10-VD-029,muc=VD,loigiai={Diện tích $S=\dfrac{1}{2}\sqrt{AB^2\cdot AC^2-(\overrightarrow{AB}\cdot\overrightarrow{AC})^2}$, $BC=\sqrt{197}$, đường cao $AH=\dfrac{2S}{BC}=\dfrac{\sqrt{449357}}{197}$.}]{Cho tam giác $ABC$ có $A(0;-1;0)$, $B(6;-2;6)$, $C(-4;-6;-3)$. Độ dài đường cao kẻ từ $A$ của tam giác bằng}
  {$\dfrac{\sqrt{449357}}{197}$.}
  {$\dfrac{\sqrt{488166}}{197}$.}{$\dfrac{\sqrt{898714}}{197}$.}{$\dfrac{\sqrt{898714}}{394}$.}

\cauTN[id=C2-B8-D10-VD-030,muc=VD,loigiai={Diện tích $S=\dfrac{1}{2}\sqrt{AB^2\cdot AC^2-(\overrightarrow{AB}\cdot\overrightarrow{AC})^2}$, $BC=\sqrt{5}$, đường cao $AH=\dfrac{2S}{BC}=\dfrac{\sqrt{1070}}{5}$.}]{Cho tam giác $ABC$ có $A(4;2;6)$, $B(-3;5;3)$, $C(-5;5;4)$. Độ dài đường cao kẻ từ $A$ của tam giác bằng}
  {$\dfrac{\sqrt{1070}}{5}$.}
  {$\dfrac{\sqrt{1095}}{5}$.}{$\dfrac{2\sqrt{535}}{5}$.}{$\dfrac{\sqrt{535}}{5}$.}

\cauTN[id=C2-B8-D10-VD-031,muc=VD,loigiai={Diện tích $S=\dfrac{1}{2}\sqrt{AB^2\cdot AC^2-(\overrightarrow{AB}\cdot\overrightarrow{AC})^2}$, $BC=\sqrt{109}$, đường cao $AH=\dfrac{2S}{BC}=\dfrac{\sqrt{128293}}{109}$.}]{Cho tam giác $ABC$ có $A(3;5;-2)$, $B(6;-2;2)$, $C(3;4;-6)$. Độ dài đường cao kẻ từ $A$ của tam giác bằng}
  {$\dfrac{\sqrt{128293}}{109}$.}
  {$\dfrac{\sqrt{140174}}{109}$.}{$\dfrac{\sqrt{256586}}{109}$.}{$\dfrac{\sqrt{256586}}{218}$.}

\cauTN[id=C2-B8-D10-VD-032,muc=VD,loigiai={Diện tích $S=\dfrac{1}{2}\sqrt{AB^2\cdot AC^2-(\overrightarrow{AB}\cdot\overrightarrow{AC})^2}$, $BC=\sqrt{41}$, đường cao $AH=\dfrac{2S}{BC}=\dfrac{2\sqrt{3034}}{41}$.}]{Cho tam giác $ABC$ có $A(-1;-5;-4)$, $B(-5;3;-5)$, $C(-3;-3;-6)$. Độ dài đường cao kẻ từ $A$ của tam giác bằng}
  {$\dfrac{2\sqrt{3034}}{41}$.}
  {$\dfrac{\sqrt{13817}}{41}$.}{$\dfrac{4\sqrt{1517}}{41}$.}{$\dfrac{2\sqrt{1517}}{41}$.}

\cauTN[id=C2-B8-D10-VD-033,muc=VD,loigiai={Diện tích $S=\dfrac{1}{2}\sqrt{AB^2\cdot AC^2-(\overrightarrow{AB}\cdot\overrightarrow{AC})^2}$, $BC=\sqrt{41}$, đường cao $AH=\dfrac{2S}{BC}=\dfrac{45\sqrt{82}}{41}$.}]{Cho tam giác $ABC$ có $A(-3;3;5)$, $B(-2;-5;-4)$, $C(0;-6;2)$. Độ dài đường cao kẻ từ $A$ của tam giác bằng}
  {$\dfrac{45\sqrt{82}}{41}$.}
  {$\dfrac{\sqrt{167731}}{41}$.}{$\dfrac{90\sqrt{41}}{41}$.}{$\dfrac{45\sqrt{41}}{41}$.}

\cauTN[id=C2-B8-D10-VD-034,muc=VD,loigiai={Diện tích $S=\dfrac{1}{2}\sqrt{AB^2\cdot AC^2-(\overrightarrow{AB}\cdot\overrightarrow{AC})^2}$, $BC=\sqrt{101}$, đường cao $AH=\dfrac{2S}{BC}=\dfrac{\sqrt{383901}}{101}$.}]{Cho tam giác $ABC$ có $A(4;4;0)$, $B(5;0;-5)$, $C(-3;-6;-6)$. Độ dài đường cao kẻ từ $A$ của tam giác bằng}
  {$\dfrac{\sqrt{383901}}{101}$.}
  {$\dfrac{\sqrt{394102}}{101}$.}{$\dfrac{\sqrt{767802}}{101}$.}{$\dfrac{\sqrt{767802}}{202}$.}

\cauTN[id=C2-B8-D10-VD-035,muc=VD,loigiai={Diện tích $S=\dfrac{1}{2}\sqrt{AB^2\cdot AC^2-(\overrightarrow{AB}\cdot\overrightarrow{AC})^2}$, $BC=\sqrt{182}$, đường cao $AH=\dfrac{2S}{BC}=\dfrac{18\sqrt{3003}}{91}$.}]{Cho tam giác $ABC$ có $A(2;-6;-1)$, $B(0;6;-5)$, $C(-6;1;6)$. Độ dài đường cao kẻ từ $A$ của tam giác bằng}
  {$\dfrac{18\sqrt{3003}}{91}$.}
  {$\dfrac{\sqrt{981253}}{91}$.}{$\dfrac{18\sqrt{6006}}{91}$.}{$\dfrac{9\sqrt{6006}}{91}$.}

\cauTN[id=C2-B8-D10-VD-036,muc=VD,loigiai={Diện tích $S=\dfrac{1}{2}\sqrt{AB^2\cdot AC^2-(\overrightarrow{AB}\cdot\overrightarrow{AC})^2}$, $BC=\sqrt{108}$, đường cao $AH=\dfrac{2S}{BC}=\dfrac{2\sqrt{2418}}{9}$.}]{Cho tam giác $ABC$ có $A(2;6;-3)$, $B(6;-4;-5)$, $C(4;-6;5)$. Độ dài đường cao kẻ từ $A$ của tam giác bằng}
  {$\dfrac{2\sqrt{2418}}{9}$.}
  {$\dfrac{\sqrt{9753}}{9}$.}{$\dfrac{4\sqrt{1209}}{9}$.}{$\dfrac{2\sqrt{1209}}{9}$.}

\cauDS[id=C2-B8-D10-VD-037,muc=VD,loigiai={Kiểm tra từng mệnh đề bằng công thức tọa độ. a) Sai. b) Đúng. c) Sai. d) Đúng.}]{Trong không gian $Oxyz$, cho $A(-2;4;5)$, $B(4;0;-4)$. Xét tính đúng sai của các mệnh đề sau:}{%
  \y{Tam giác $ABC$ có $A(2;1;2)$, $B(3;3;4)$, $C(3;-1;4)$ có chu vi bằng $11$.}{s}
  \y{Điểm $M$ thuộc trục $Oy$ cách đều $A$ và $B$ có tọa độ $(0;\dfrac{13}{8};0)$.}{d}
  \y{Tam giác $ABC$ có $A(-2;4;5)$, $B(4;0;-4)$, $C(1;2;-3)$ thì đường cao kẻ từ $A$ bằng $\dfrac{\sqrt{186}}{2}$.}{s}
  \y{Tam giác $ABC$ có $A(-2;4;5)$, $B(4;0;-4)$, $C(1;2;-3)$ thì đường cao kẻ từ $A$ bằng $\dfrac{\sqrt{182}}{2}$.}{d}}

\cauDS[id=C2-B8-D10-VD-038,muc=VD,loigiai={Kiểm tra từng mệnh đề bằng công thức tọa độ. a) Đúng. b) Sai. c) Đúng. d) Sai.}]{Trong không gian $Oxyz$, cho $A(5;6;3)$, $B(0;-5;3)$, $C(-3;0;3)$. Xét tính đúng sai của các mệnh đề sau:}{%
  \y{Điểm $M$ thuộc trục $Ox$ cách đều $A$ và $B$ có tọa độ $(\dfrac{18}{5};0;0)$.}{d}
  \y{Diện tích tam giác $ABC$ bằng $29\sqrt{2}$.}{s}
  \y{Diện tích tam giác $ABC$ bằng $29$.}{d}
  \y{Tam giác $ABC$ có $A(5;6;3)$, $B(0;-5;3)$, $C(-3;0;3)$ thì đường cao kẻ từ $A$ bằng $\dfrac{58\sqrt{17}}{17}$.}{s}}

\cauDS[id=C2-B8-D10-VD-039,muc=VD,loigiai={Kiểm tra từng mệnh đề bằng công thức tọa độ. a) Sai. b) Sai. c) Đúng. d) Đúng.}]{Trong không gian $Oxyz$, cho $A(-5;0;-2)$, $B(-1;0;-2)$, $C(-4;0;-4)$. Xét tính đúng sai của các mệnh đề sau:}{%
  \y{Tam giác $ABC$ có $A(-5;0;-2)$, $B(-1;0;-2)$, $C(-4;0;-4)$ thì đường cao kẻ từ $A$ bằng $\dfrac{4\sqrt{26}}{13}$.}{s}
  \y{Tam giác $ABC$ có $A(-5;0;-2)$, $B(-1;0;-2)$, $C(-4;0;-4)$ thì đường cao kẻ từ $A$ bằng $\dfrac{8\sqrt{26}}{13}$.}{s}
  \y{Diện tích tam giác $ABC$ bằng $4$.}{d}
  \y{Tam giác $ABC$ có $A(1;1;3)$, $B(2;3;5)$, $C(2;-1;5)$ có chu vi bằng $10$.}{d}}

\cauDS[id=C2-B8-D10-VD-040,muc=VD,loigiai={Kiểm tra từng mệnh đề bằng công thức tọa độ. a) Đúng. b) Sai. c) Đúng. d) Sai.}]{Trong không gian $Oxyz$, cho $A(-6;2;-1)$, $B(-1;0;4)$, $C(6;-4;-1)$. Xét tính đúng sai của các mệnh đề sau:}{%
  \y{Diện tích tam giác $ABC$ bằng $9\sqrt{14}$.}{d}
  \y{Tam giác $ABC$ có $A(-1;-3;3)$, $B(-1;-3;8)$, $C(3;1;10)$ có chu vi bằng $19$.}{s}
  \y{Tam giác $ABC$ có $A(-1;-1;3)$, $B(-3;1;4)$, $C(1;5;6)$ có chu vi bằng $16$.}{d}
  \y{Tam giác $ABC$ có $A(-6;2;-1)$, $B(-1;0;4)$, $C(6;-4;-1)$ thì đường cao kẻ từ $A$ bằng $\dfrac{3\sqrt{70}}{5}$.}{s}}

\cauDS[id=C2-B8-D10-VD-041,muc=VD,loigiai={Kiểm tra từng mệnh đề bằng công thức tọa độ. a) Sai. b) Đúng. c) Sai. d) Sai.}]{Trong không gian $Oxyz$, cho $A(5;4;0)$, $B(-3;4;-1)$, $C(3;5;5)$. Xét tính đúng sai của các mệnh đề sau:}{%
  \y{Diện tích tam giác $ABC$ bằng $\sqrt{1829}$.}{s}
  \y{Tam giác $ABC$ có $A(0;-2;2)$, $B(1;-2;2)$, $C(5;-2;2)$ có chu vi bằng $10$.}{d}
  \y{Diện tích tam giác $ABC$ bằng $\dfrac{\sqrt{3658}}{2}$.}{s}
  \y{Tam giác $ABC$ có $A(5;4;0)$, $B(-3;4;-1)$, $C(3;5;5)$ thì đường cao kẻ từ $A$ bằng $\dfrac{\sqrt{267034}}{73}$.}{s}}

\cauDS[id=C2-B8-D10-VD-042,muc=VD,loigiai={Kiểm tra từng mệnh đề bằng công thức tọa độ. a) Đúng. b) Đúng. c) Sai. d) Đúng.}]{Trong không gian $Oxyz$, cho $A(6;0;5)$, $B(0;-4;5)$, $C(-6;-1;6)$. Xét tính đúng sai của các mệnh đề sau:}{%
  \y{Tam giác $ABC$ có $A(6;0;5)$, $B(0;-4;5)$, $C(-6;-1;6)$ thì đường cao kẻ từ $A$ bằng $\dfrac{2\sqrt{5221}}{23}$.}{d}
  \y{Diện tích tam giác $ABC$ bằng $\sqrt{454}$.}{d}
  \y{Tam giác $ABC$ có $A(6;0;5)$, $B(0;-4;5)$, $C(-6;-1;6)$ thì đường cao kẻ từ $A$ bằng $\dfrac{\sqrt{10442}}{23}$.}{s}
  \y{Điểm $M$ thuộc trục $Oy$ cách đều $A$ và $B$ có tọa độ $(0;\dfrac{5}{2};0)$.}{d}}

\cauDS[id=C2-B8-D10-VD-043,muc=VD,loigiai={Kiểm tra từng mệnh đề bằng công thức tọa độ. a) Sai. b) Đúng. c) Sai. d) Sai.}]{Trong không gian $Oxyz$, cho $A(-1;-1;-4)$, $B(-6;2;2)$. Xét tính đúng sai của các mệnh đề sau:}{%
  \y{Tam giác $ABC$ có $A(-2;3;3)$, $B(2;7;10)$, $C(-2;3;8)$ có chu vi bằng $22$.}{s}
  \y{Điểm $M$ thuộc trục $Oz$ cách đều $A$ và $B$ có tọa độ $(0;0;\dfrac{13}{6})$.}{d}
  \y{Tam giác $ABC$ có $A(-1;-1;-4)$, $B(-6;2;2)$, $C(4;-4;0)$ thì đường cao kẻ từ $A$ bằng $\dfrac{\sqrt{1239}}{7}$.}{s}
  \y{Tam giác $ABC$ có $A(-1;-1;-4)$, $B(-6;2;2)$, $C(4;-4;0)$ thì đường cao kẻ từ $A$ bằng $\dfrac{2\sqrt{595}}{7}$.}{s}}

\cauDS[id=C2-B8-D10-VD-044,muc=VD,loigiai={Kiểm tra từng mệnh đề bằng công thức tọa độ. a) Đúng. b) Sai. c) Đúng. d) Sai.}]{Trong không gian $Oxyz$, cho $A(4;0;2)$, $B(6;2;-2)$. Xét tính đúng sai của các mệnh đề sau:}{%
  \y{Tam giác $ABC$ có $A(4;0;2)$, $B(6;2;-2)$, $C(4;5;4)$ thì đường cao kẻ từ $A$ bằng $\dfrac{2\sqrt{173}}{7}$.}{d}
  \y{Tam giác $ABC$ có $A(-3;2;3)$, $B(-1;4;2)$, $C(-1;4;4)$ có chu vi bằng $10$.}{s}
  \y{Điểm $M$ thuộc trục $Oz$ cách đều $A$ và $B$ có tọa độ $(0;0;-3)$.}{d}
  \y{Tam giác $ABC$ có $A(4;0;2)$, $B(6;2;-2)$, $C(4;5;4)$ thì đường cao kẻ từ $A$ bằng $\dfrac{2\sqrt{346}}{7}$.}{s}}

\cauDS[id=C2-B8-D10-VD-045,muc=VD,loigiai={Kiểm tra từng mệnh đề bằng công thức tọa độ. a) Sai. b) Đúng. c) Đúng. d) Đúng.}]{Trong không gian $Oxyz$, cho $A(-1;6;-6)$, $B(5;-3;5)$. Xét tính đúng sai của các mệnh đề sau:}{%
  \y{Tam giác $ABC$ có $A(-1;6;-6)$, $B(5;-3;5)$, $C(6;-6;-6)$ thì đường cao kẻ từ $A$ bằng $\dfrac{2\sqrt{1534927}}{131}$.}{s}
  \y{Tam giác $ABC$ có $A(2;-1;3)$, $B(4;1;4)$, $C(0;1;4)$ có chu vi bằng $10$.}{d}
  \y{Điểm $M$ thuộc trục $Ox$ cách đều $A$ và $B$ có tọa độ $(-\dfrac{7}{6};0;0)$.}{d}
  \y{Tam giác $ABC$ có $A(-1;6;-6)$, $B(5;-3;5)$, $C(6;-6;-6)$ thì đường cao kẻ từ $A$ bằng $\dfrac{\sqrt{3069854}}{131}$.}{d}}

\cauDS[id=C2-B8-D10-VD-046,muc=VD,loigiai={Kiểm tra từng mệnh đề bằng công thức tọa độ. a) Đúng. b) Đúng. c) Sai. d) Sai.}]{Trong không gian $Oxyz$, cho $A(-4;6;-5)$, $B(-5;6;4)$. Xét tính đúng sai của các mệnh đề sau:}{%
  \y{Tam giác $ABC$ có $A(-3;-1;2)$, $B(-3;0;2)$, $C(-3;4;2)$ có chu vi bằng $10$.}{d}
  \y{Tam giác $ABC$ có $A(-4;6;-5)$, $B(-5;6;4)$, $C(0;-3;-2)$ thì đường cao kẻ từ $A$ bằng $\dfrac{3\sqrt{128794}}{142}$.}{d}
  \y{Điểm $M$ thuộc trục $Oz$ cách đều $A$ và $B$ có tọa độ $(0;0;2)$.}{s}
  \y{Điểm $M$ thuộc trục $Ox$ cách đều $A$ và $B$ có tọa độ $(-2;0;0)$.}{s}}

\cauDS[id=C2-B8-D10-VD-047,muc=VD,loigiai={Kiểm tra từng mệnh đề bằng công thức tọa độ. a) Đúng. b) Sai. c) Sai. d) Sai.}]{Trong không gian $Oxyz$, cho $A(2;1;-2)$, $B(-6;5;3)$, $C(-1;1;-4)$. Xét tính đúng sai của các mệnh đề sau:}{%
  \y{Tam giác $ABC$ có $A(-3;-2;0)$, $B(-3;3;0)$, $C(-3;-1;0)$ có chu vi bằng $10$.}{d}
  \y{Tam giác $ABC$ có $A(2;1;-2)$, $B(-6;5;3)$, $C(-1;1;-4)$ thì đường cao kẻ từ $A$ bằng $\dfrac{\sqrt{12590}}{30}$.}{s}
  \y{Diện tích tam giác $ABC$ bằng $\dfrac{\sqrt{1173}}{2}$.}{s}
  \y{Tam giác $ABC$ có $A(-2;2;2)$, $B(-4;4;3)$, $C(0;4;3)$ có chu vi bằng $11$.}{s}}

\cauDS[id=C2-B8-D10-VD-048,muc=VD,loigiai={Kiểm tra từng mệnh đề bằng công thức tọa độ. a) Sai. b) Sai. c) Sai. d) Đúng.}]{Trong không gian $Oxyz$, cho $A(-4;6;-6)$, $B(5;6;-1)$, $C(2;4;-4)$. Xét tính đúng sai của các mệnh đề sau:}{%
  \y{Điểm $M$ thuộc trục $Ox$ cách đều $A$ và $B$ có tọa độ $(-\dfrac{13}{9};0;2)$.}{s}
  \y{Diện tích tam giác $ABC$ bằng $2\sqrt{142}$.}{s}
  \y{Tam giác $ABC$ có $A(-4;6;-6)$, $B(5;6;-1)$, $C(2;4;-4)$ thì đường cao kẻ từ $A$ bằng $\dfrac{\sqrt{3245}}{11}$.}{s}
  \y{Tam giác $ABC$ có $A(-3;-3;3)$, $B(-1;-1;4)$, $C(-1;-1;2)$ có chu vi bằng $8$.}{d}}

\cauTLN[id=C2-B8-D10-VD-049,muc=VD,loigiai={Gọi $M(x;y;0)$. Từ $MA^2=MB^2$ và $MA^2=MC^2$ ta được hệ hai phương trình bậc nhất ẩn $x,y$, giải được $x=-1$, $y=-3$. Vậy $M(-1;-3;0)$.}]{Cho $A(0;-1;2)$, $B(1;-2;2)$, $C(1;-4;2)$. Điểm $M$ thuộc mặt phẳng $(Oxy)$ cách đều ba điểm $A,B,C$. Tính $x_M+y_M$.}{-4}

\cauTLN[id=C2-B8-D10-VD-050,muc=VD,loigiai={Gọi $M(x;y;0)$. Từ $MA^2=MB^2$ và $MA^2=MC^2$ ta được hệ hai phương trình bậc nhất ẩn $x,y$, giải được $x=1$, $y=-1$. Vậy $M(1;-1;0)$.}]{Cho $A(2;-3;2)$, $B(3;1;1)$, $C(3;-2;2)$. Điểm $M$ thuộc mặt phẳng $(Oxy)$ cách đều ba điểm $A,B,C$. Tính $x_M+y_M$.}{0}

\cauTLN[id=C2-B8-D10-VD-051,muc=VD,loigiai={Gọi $M(x;y;0)$. Từ $MA^2=MB^2$ và $MA^2=MC^2$ ta được hệ hai phương trình bậc nhất ẩn $x,y$, giải được $x=2$, $y=-1$. Vậy $M(2;-1;0)$.}]{Cho $A(4;2;6)$, $B(8;2;2)$, $C(5;1;6)$. Điểm $M$ thuộc mặt phẳng $(Oxy)$ cách đều ba điểm $A,B,C$. Tính $x_M+y_M$.}{1}

\cauTLN[id=C2-B8-D10-VD-052,muc=VD,loigiai={Gọi $M(x;y;0)$. Từ $MA^2=MB^2$ và $MA^2=MC^2$ ta được hệ hai phương trình bậc nhất ẩn $x,y$, giải được $x=1$, $y=-2$. Vậy $M(1;-2;0)$.}]{Cho $A(1;2;3)$, $B(5;1;0)$, $C(1;3;0)$. Điểm $M$ thuộc mặt phẳng $(Oxy)$ cách đều ba điểm $A,B,C$. Tính $x_M+y_M$.}{-1}

\cauTLN[id=C2-B8-D10-VD-053,muc=VD,loigiai={Gọi $M(x;y;0)$. Từ $MA^2=MB^2$ và $MA^2=MC^2$ ta được hệ hai phương trình bậc nhất ẩn $x,y$, giải được $x=0$, $y=2$. Vậy $M(0;2;0)$.}]{Cho $A(3;6;0)$, $B(4;2;3)$, $C(4;5;0)$. Điểm $M$ thuộc mặt phẳng $(Oxy)$ cách đều ba điểm $A,B,C$. Tính $x_M+y_M$.}{2}

\cauTLN[id=C2-B8-D10-VD-054,muc=VD,loigiai={Gọi $M(x;y;0)$. Từ $MA^2=MB^2$ và $MA^2=MC^2$ ta được hệ hai phương trình bậc nhất ẩn $x,y$, giải được $x=-1$, $y=0$. Vậy $M(-1;0;0)$.}]{Cho $A(1;2;-1)$, $B(1;1;2)$, $C(-3;2;1)$. Điểm $M$ thuộc mặt phẳng $(Oxy)$ cách đều ba điểm $A,B,C$. Tính $x_M+y_M$.}{-1}

\cauTLN[id=C2-B8-D10-VD-055,muc=VD,loigiai={$AB=3$, $AC=3$, $BC=4$ nên chu vi bằng $3+3+4=10$.}]{Tam giác $ABC$ có $A(3;-3;0)$, $B(4;-5;2)$, $C(4;-1;2)$. Tính chu vi tam giác $ABC$.}{10}

\cauTLN[id=C2-B8-D10-VD-056,muc=VD,loigiai={$AB=5$, $AC=1$, $BC=4$ nên chu vi bằng $5+1+4=10$.}]{Tam giác $ABC$ có $A(-3;-3;1)$, $B(-3;-3;6)$, $C(-3;-3;2)$. Tính chu vi tam giác $ABC$.}{10}

\cauTLN[id=C2-B8-D10-VD-057,muc=VD,loigiai={$AB=7$, $AC=3$, $BC=6$ nên chu vi bằng $7+3+6=16$.}]{Tam giác $ABC$ có $A(-2;-2;-3)$, $B(0;4;0)$, $C(-4;0;-2)$. Tính chu vi tam giác $ABC$.}{16}

\cauTLN[id=C2-B8-D10-VD-058,muc=VD,loigiai={$AB=5$, $AC=1$, $BC=4$ nên chu vi bằng $5+1+4=10$.}]{Tam giác $ABC$ có $A(-2;2;3)$, $B(3;2;3)$, $C(-1;2;3)$. Tính chu vi tam giác $ABC$.}{10}

\cauTLN[id=C2-B8-D10-VD-059,muc=VD,loigiai={$AB=5$, $AC=1$, $BC=4$ nên chu vi bằng $5+1+4=10$.}]{Tam giác $ABC$ có $A(-1;-1;-3)$, $B(-1;-1;2)$, $C(-1;-1;-2)$. Tính chu vi tam giác $ABC$.}{10}

\cauTLN[id=C2-B8-D10-VD-060,muc=VD,loigiai={$AB=3$, $AC=7$, $BC=6$ nên chu vi bằng $3+7+6=16$.}]{Tam giác $ABC$ có $A(2;-1;2)$, $B(0;1;3)$, $C(4;5;5)$. Tính chu vi tam giác $ABC$.}{16}
